% 本文件由 LaTeX 排版 Agent 根据有效 translation.json 自动生成。
% author=Jerome; work_id=2708; title=论著名人物

\chapter{论著名人物}

% structure: p0001 source=h1 target=omitted reason=page-title-already-rendered-by-parent

你曾催促我，\Person{德克斯特}，要效法\Person{特兰奎卢斯}的榜样，对教会作家作系统的记述，并为我们的人做他对付外邦人中的著名文人所做的事，即，简要地向你介绍所有那些从主的受难直到\Person{德奥多西}皇帝第十四年间，曾发表过关于圣经的任何值得纪念的著作的人。类似的工作，希腊人中逍遥派的\Person{赫米普斯}、\Person{卡里斯提的安提戈诺斯}、博学的\Person{萨提罗斯}，以及所有中最博学的音乐家\Person{阿里斯托克塞努斯}已经完成；拉丁人中，\Person{瓦罗}、\Person{桑特拉}、\Person{内波斯}、\Person{希吉努斯}，以及你想以其榜样激励我们的\Person{特兰奎卢斯}也已完成。然而，他们的情况与我的情况并不相同，因为他们打开古代的历史和编年史，仿佛从某个大草地上采集，至少可以为他们的作品编织一个小小的花环。至于我，该做什么呢？我没有前人，如俗语所说，我有最糟糕的老师，即我自己，但我必须承认，\Person{尤西比乌·庞菲卢斯}在他的十卷\WorkTitle{教会史}中给了我极大的帮助，而我们将要述及的许多人的著作本身也常常为其作者的生平年代作证。因此，我祈求主耶稣，使你那站在罗马雄辩术顶峰的\Person{西塞罗}不耻于做的，即在他的\WorkTitle{布鲁图斯}中编列拉丁演说家的名录，我也能在对教会作家的叙述中完成，并以一种配得上你所提出的劝勉的方式完成。但是，如果在这卷书中，我偶然遗漏了那些仍在写作的人中的某一位，他们应当将此归咎于自己，而非归咎于我，因为在我未曾读过的那些人中，我首先无法知道那些隐藏自己著作的人，其次，那些也许为他人所熟知的事，在我这个偏僻的地球角落，可能完全不为我所知。但是，当他们以其著作而闻名时，他们不会因我们未曾提及而过分悲痛。让反对基督的狂怒者——\Person{克耳苏斯}、\Person{波斐利}、\Person{犹利安}——以及他们的追随者，那些以为教会从未有过哲学家、演说家或学者的人，知道有多少和什么样的人奠定了教会的基础、建造并装饰了它，并停止指责我们的信仰如此粗俗简单，而应认识到他们自己的无知。\par

然而，他们的情况与我的情况并不相同，因为他们打开古代的历史和编年史，仿佛从某个大草地上采集，至少可以为他们的作品编织一个小小的花环。至于我，该做什么呢？我没有前人，如俗语所说，我有最糟糕的老师，即我自己，但我必须承认，\Person{尤西比乌·庞菲卢斯}在他的十卷\WorkTitle{教会史}中给了我极大的帮助，而我们将要述及的许多人的著作本身也常常为其作者的生平年代作证。因此，我祈求主耶稣，使你那站在罗马雄辩术顶峰的\Person{西塞罗}不耻于做的，即在他的\WorkTitle{布鲁图斯}中编列拉丁演说家的名录，我也能在对教会作家的叙述中完成，并以一种配得上你所提出的劝勉的方式完成。但是，如果在这卷书中，我偶然遗漏了那些仍在写作的人中的某一位，他们应当将此归咎于自己，而非归咎于我，因为在我未曾读过的那些人中，我首先无法知道那些隐藏自己著作的人，其次，那些也许为他人所熟知的事，在我这个偏僻的地球角落，可能完全不为我所知。但是，当他们以其著作而闻名时，他们不会因我们未曾提及而过分悲痛。让反对基督的狂怒者——\Person{克耳苏斯}、\Person{波斐利}、\Person{犹利安}——以及他们的追随者，那些以为教会从未有过哲学家、演说家或学者的人，知道有多少和什么样的人奠定了教会的基础、建造并装饰了它，并停止指责我们的信仰如此粗俗简单，而应认识到他们自己的无知。\par

% structure: p0004 source=h2 target=section reason=relative-heading-rank; base=h2

\section{第一章}

\Person{西满伯多禄}，若望之子，来自\Place{加里肋亚}省的\Place{贝特赛达}村，是宗徒\Person{安德肋}的兄弟，他自己是宗徒之长，在担任\Place{安提约基雅}教会的主教并向散居的犹太信徒——那些受割礼的信徒——在\Place{本都}、\Place{迦拉达}、\Place{卡帕多细亚}、\Place{亚细亚}和\Place{彼提尼雅}宣讲之后，在\Person{克劳狄}第二年前往\Place{罗马}推翻\Person{西满魔术师}，并在那里担任主教职位二十五年，直到\Person{尼禄}的最后一年，即第十四年。他因被钉十字架而接受了殉道的荣冠，头朝下、脚朝天，声称他不配以与他的主相同的方式被钉死。他写了两封被称为公函的书信，第二封由于风格与第一封不同，被许多人认为并非他所写。此外，那部根据他的门徒和译者\Person{马尔谷}所写的福音也归于他。另一方面，那些著作中，一本称为他的\WorkTitle{宗徒大事录}，另一本他的\WorkTitle{福音}，第三本他的\WorkTitle{宣讲}，第四本他的\WorkTitle{默示录}，第五本他的\Quote{审判}，都被视为伪经而予以拒绝。\par

他被葬于\Place{罗马}\Place{梵蒂冈}靠近凯旋大道的墓地，受到全世界的敬仰。\par

% structure: p0007 source=h2 target=section reason=relative-heading-rank; base=h2

\section{第二章}

被称为主的兄弟的\Person{雅各伯}，别号义者，如有些人所认为，是\Person{若瑟}由另一妻子所生，但依我看来，是\Person{玛利亚}——我们的主的母亲的姐妹——的儿子，\Person{若望}在他的书中提到过她。在主的受难之后，他立即被宗徒们任命为\Place{耶路撒冷}的主教，写了一封单独的书信，这封信被列入七封公函之一，甚至这封信也被一些人声称是由别人以他的名义发表，而后随着时间的推移逐渐获得了权威。接近宗徒时代的\Person{赫格西普斯}，在他的\WorkTitle{注释}第五卷中论及\Person{雅各伯}时说：\par

\Quote{宗徒之后，主的兄弟\Person{雅各伯}，别号义者，被立为\Place{耶路撒冷}教会的首领。确实有许多人称为\Person{雅各伯}。这位从母胎中就是圣的。他不喝酒和烈酒，不吃肉，从不剃头或用油膏抹自己，也不洗澡。他独享进入至圣所的特权，因为他确实不穿羊毛衣服，而穿亚麻衣，独自进入圣殿，为百姓祈祷，以至于他的膝盖据说已变得如骆驼膝盖般坚硬。}\par

他还说了许多其他事，多到不可胜数。若瑟夫在他的\WorkTitle{犹太古史}第二十卷，以及克莱孟在他的\WorkTitle{纲要}第七卷中也提到：斐斯托统治\Place{犹太}去世后，尼禄派遣阿尔比努斯继任。在他到达行省之前，大司铎阿纳尼亚——司铎阶层阿纳斯的年轻儿子——趁无政府状态，召集公议会，公开试图强迫雅各伯否认基督是天主子。当他拒绝时，阿纳尼亚下令用石头砸死他。从圣殿的屋顶被推下，腿骨摔断，但还半活，他向天举起双手说，\BibleQuote{主啊，宽恕他们，因为他们不知道自己在做什么。}接着，被漂布者用来拧干衣服的棒子击中头部——他便死了。这位若瑟夫记载了一个传说：这位雅各伯在民众中享有如此大的圣德和声望，以至于人们相信耶路撒冷的陷落是因他的死亡所致。正是他，宗徒保禄在写给迦拉达人的书信中写道：\BibleQuote{除了主的兄弟雅各伯，我没有看见别的宗徒}，不久之后，\BibleBook{宗徒}{大事录}也为此事作证。那部称为\WorkTitle{希伯来人福音}的福音——我最近将其译成希腊文和拉丁文，奥力振也经常使用它——在记载救主复活之后说：\Quote{但是主在把殓布交给司铎的仆人之后，显现给雅各伯（因为雅各伯曾发誓，从他喝主的杯的那一刻起，直到看见他从睡眠中复活，他就不吃饼）}稍后又说：\Quote{\InnerQuote{拿张桌子，拿饼来}，主说。}接着立即补充说：\Quote{他拿来了饼，祝福了，掰开，递给义人雅各伯，并对他说：\InnerQuote{我的兄弟，吃你的饼吧，因为人子已从睡眠中复活了。}}于是他治理耶路撒冷教会三十年，即直到尼禄第七年，并被葬在他被推下的圣殿附近。他的墓碑及其铭文为人熟知，直到提托围城和哈德良统治结束。我们的一些作者认为他被葬在橄榄山，但他们弄错了。\par

% structure: p0011 source=h2 target=section reason=relative-heading-rank; base=h2

\section{第三章}

\Person{玛窦}，也称为\Person{肋未}，宗徒、昔为税吏，撰写了一部基督的福音，最初在\Place{犹太}以希伯来文出版，为那些受割礼而信的人；但后来被译成希腊文，不过译者为谁不确定。希伯来原文至今仍保存在\Place{凯撒勒雅}的图书馆中，那是\Person{庞斐禄}勤奋收集的。我也有机会请叙利亚城市\Place{贝洛雅}的纳匝肋人将那卷书描述给我，因为他们使用它。在这卷书中需注意，每当福音作者——无论以他自己的身份，还是以我们的主救主的口吻——引用旧约的见证时，他都不遵从七十贤士译本译者的权威，而是跟随希伯来原文。因此这两个形式存在：\BibleQuote{我从埃及召回了我的儿子}和\BibleQuote{他将被称为纳匝肋人。}\par

% structure: p0013 source=h2 target=section reason=relative-heading-rank; base=h2

\section{第四章}

\Person{犹达}，\Person{雅各伯}的兄弟，留下了一封简短的书信，这封信被列入七封公函之中；因为他在其中引用了伪经\WorkTitle{哈诺客书}，所以被许多人拒绝。然而，因年代久远和持续使用，它已获得权威，并被列入圣经。\par

% structure: p0015 source=h2 target=section reason=relative-heading-rank; base=h2

\section{第五章}

保禄，原名扫禄，是不在十二宗徒之列的宗徒，属本雅明支派，生于犹太的吉斯卡拉镇。该镇被罗马人占领后，他随父母迁往基里基雅的塔尔索。被父母送往耶路撒冷学习法律，受教于加玛里尔，一位路加提及的博学之士。但他曾目睹殉道者斯德望的死亡，并从圣殿的大司祭取得信函，为迫害信奉基督的人，前往大马士革。在那里，因启示而被迫归信，正如\BibleBook{}{宗徒大事录}所载，他从迫害者转变为特选的器皿。由于塞浦路斯总督色尔爵·保禄首先相信他的宣讲，他便取了此人的名字，因为他使他归顺于对基督的信仰。他与巴尔纳伯会合后，遍历许多城市，回到耶路撒冷，由伯多禄、雅各伯和若望祝圣为外邦人的宗徒。由于\BibleBook{}{宗徒大事录}已详细记述了他的生平，我只说这一点：吾主受难后第二十五年，即尼禄第二年，在犹太总督斐斯托接替斐理斯时，他被捆绑解送罗马，在自由拘留所住了两年，天天与犹太人辩论关于基督的来临。应当说明，在第一次辩护时，尼禄的权势尚未稳固，他的邪恶也未如史书所述那般爆发，保禄被尼禄释放，为使基督的福音也能在西方宣讲。正如他在\BibleBook{弟茂德}{后书}中所写的，当时他即将受死，在锁链中口授此书信：\BibleQuote{在我初次辩护时，没有人支持我，反而都离弃了我；愿这罪不归于他们。但主站在我身旁，坚固了我，好使我藉着我完全宣讲这信息，并使所有外邦人都能听见；我也从狮子的口中被救出来。}——显然以狮子指尼禄，因其残暴。接着他说道：\BibleQuote{主救我从狮子的口中出来}，稍后又\BibleQuote{主救我脱离一切凶恶之事，并救我进入他的天国}，因为他内心感到自己的殉道临近了，在同一书信中他宣布：\BibleQuote{因为我已被奠祭，我离世的时期到了。}然后，在尼禄第十四年，与伯多禄同日，他为基督在罗马被斩首，葬于奥斯提亚道，时在吾主受难后第二十七年。他写给七个教会九封书信：\WorkTitle{致罗马人}一封、\WorkTitle{致格林多人}两封、\WorkTitle{致迦拉达人}一封、\WorkTitle{致厄弗所人}一封、\WorkTitle{致斐理伯人}一封、\WorkTitle{致哥罗森人}一封、\WorkTitle{致得撒洛尼人}两封；此外还有给其门徒的：\WorkTitle{致弟茂德}两封、\WorkTitle{致弟铎}一封、\WorkTitle{致费肋孟}一封。那称为\WorkTitle{致希伯来人}的书信不被认为是他的，因其风格和语言与其他书信不同；但据戴尔都良之说，它被认为是巴尔纳伯的作品；或据他人之说，是路加福音作者或后来罗马教会主教克肋孟的作品，他们声称将保禄的思想用自己的语言整理润饰；当然，因为保禄是写给希伯来人，且在他们中不受欢迎，他可能因此省略了自己的署名。他是希伯来人，用希伯来文写作，即他的母语且极其流畅，而用希伯来文优美写成的作品被更优美地译成希腊文，这就是为何它似乎不同于保禄的其他书信。也有人读到一封\WorkTitle{致劳狄刻雅人}的书信，但所有人都拒绝它。\par

% structure: p0017 source=h2 target=section reason=relative-heading-rank; base=h2

\section{第六章}

巴尔纳伯，塞浦路斯人，又名若瑟肋未人，与保禄一同被祝圣为外邦人的宗徒，写了一封\Emph{书信}，有益于教会的建立，被列为次经著作。后来因门徒若望（又名马尔谷）而与保禄分离，但并未因此懈怠他所领受的宣讲福音的使命。\par

% structure: p0019 source=h2 target=section reason=relative-heading-rank; base=h2

\section{第七章}

路加，安提约基雅的医师，如其著作所示，精通希腊文。他是保禄宗徒的追随者，伴其所有旅程，写了一部\BibleBook{路加}{福音}。关于此书，同一保禄说：\BibleQuote{我们打发一位弟兄与他同去，这位弟兄在福音上因所有教会而受称赞。}又对哥罗森人说：\BibleQuote{亲爱的医生路加问候你们。}对弟茂德说：\BibleQuote{只有路加与我在一起。}他还写了另一部极好的著作，命名为\BibleBook{}{宗徒大事录}，一部历史，延伸至保禄在罗马居住的第二年，即尼禄第四年，由此我们得知该书是在同一城市写成的。因此，我们将\WorkTitle{保禄与德克拉行传}以及关于被他施洗的狮子的所有寓言视为次经著作。因为，怎么可能这位宗徒在其他事务中形影不离的同伴，唯独对此事一无所知呢？此外，戴尔都良，生活于接近那些时代，提到亚细亚的一位长老，保禄宗徒的追随者，被若望断定是此书的作者，他承认自己这样做是出于对保禄的爱，因而辞去了长老职务。有些人认为，每当保禄在他的书信中说\BibleQuote{按我的福音}，他指的是路加的书，而路加不仅从保禄宗徒学习了福音历史（保禄并未在肉身中与主同在），也从其他宗徒那里学习。他自己在著作开头也声明说：\BibleQuote{正如那些从起初亲眼看见并为话语的仆役的人，传给我们一样。}因此，他按所听到的写了福音，但按自己所看到的编写了\BibleBook{}{宗徒大事录}。他被葬于君士坦丁堡，在君士坦提乌斯第二十年，他的骸骨与安德肋宗徒的遗骸一同被迁往该城。\par

% structure: p0021 source=h2 target=section reason=relative-heading-rank; base=h2

\section{第八章}

玛尔谷，伯多禄的门徒和翻译者，应罗马弟兄们的请求，根据他所听到的伯多禄的讲述，撰写了一部简短的福音。伯多禄听说后，予以批准，并以其权威将该福音公布于各教会供人诵读，正如克莱孟在其\WorkTitle{要义}第六卷中所述，以及希拉波利斯的主教帕皮亚所记载。伯多禄也在他的第一封书信中提到了这位玛尔谷，以巴比伦之名象征性地指代罗马，说：\Quote{在巴比伦同蒙拣选的教会问候你们，我儿玛尔谷也问候你们。}于是，他带着自己所撰写的福音，前往埃及，在亚历山大首先宣讲基督，建立了一个在教义和生活节制方面令人钦佩的教会，以至于他迫使所有基督的追随者效法他的榜样。犹太人中最博学的斐罗，看到亚历山大初期的教会仍带有犹太色彩，便撰写了一部关于他们生活方式的著作，作为对其民族的光荣，记述了如同路加所说，耶路撒冷的信徒凡物公用，他记载在学识渊博的玛尔谷领导下，在亚历山大所见到的同样情况。玛尔谷在尼禄在位第八年去世，葬于亚历山大，由安尼雅诺继任。\par

% structure: p0023 source=h2 target=section reason=relative-heading-rank; base=h2

\section{第九章}

若望，耶稣所爱的宗徒，载伯德的儿子，雅各伯的兄弟——这位雅各伯在吾主受难后被黑落德斩首——是所有福音作者中最晚撰写一部\WorkTitle{福音}的。他应亚细亚各主教的要求，为反对则林特及其他异端者，尤其是当时正在兴起的厄彼翁派（他们声称基督在玛利亚之前并不存在）的教义而作。为此，他不得不维护基督的神圣诞生。但关于这部作品，据说还有另一个原因：当他阅读了玛窦、玛尔谷和路加的福音后，他固然认可了历史的实质内容，并宣称他们所说的是真实的，但他们只记述了一年的历史，即随若翰被囚禁之后、他殉道的那一年。因此，他略过这一年（这些作者已叙述了这一年的事件），记述了若翰被囚禁之前更早时期的事件，以便那些勤于阅读四部福音书的人能够明瞭。这也消除了若望与其他福音作者之间似乎存在的差异。他还撰写了一封\WorkTitle{书信}，开头如下：\Quote{论到那从起初原有的生命之圣言，就是我们听见过、亲眼看见过、注视过、亲手摸过的。}这封书信受到所有关注教会或学识之人的珍视。另外两封书信，第一封开头为\Quote{长老致蒙拣选的夫人和她的儿女}，另一封为\Quote{长老致亲爱的加约，就是我在真理中所爱的}，据说是长老若望的作品，至今在厄弗所仍另有一座他的坟墓，尽管有些人认为这是同一位福音作者若望的两处纪念地。我们将在适当时候，在谈到他的门徒帕皮亚时讨论此事。在图密善继尼禄之后发动第二次迫害的第十四年，若望被放逐到帕特摩岛，并撰写了\WorkTitle{默示录}，后来儒斯定·玛尔蒂尔和依肋内曾为其作注释。但图密善被处死，因其过度残暴，元老院废除了他的法令，若望在佩蒂纳克斯统治时期回到厄弗所，一直居住到图拉真皇帝时代，在整个亚细亚建立并修建了教会，因年老力衰，在吾主受难后第六十八年去世，葬于同一城市附近。\par

% structure: p0025 source=h2 target=section reason=relative-heading-rank; base=h2

\section{第十章}

赫尔玛，宗徒保禄在致罗马人书中提到他说：\Quote{请问候弗肋贡、赫尔默斯、帕特洛巴、赫尔玛和与他们在一起的弟兄们。}他被认为是称为\WorkTitle{牧者}那部书的作者，该书在希腊的某些教会中也公开宣读。它确实是一部有益的书，许多古代作者引用它作为权威，但在拉丁教会中几乎不为人知。\par

% structure: p0027 source=h2 target=section reason=relative-heading-rank; base=h2

\section{第十一章}

犹太人\Person{斐洛}，亚历山大里亚人，属于司祭阶级，我们将其列入教会作家之列，理由是他撰写了关于马尔谷福音者在亚历山大里亚的第一个教会的书，并在其中赞扬我们，声明不仅他们在那里，而且在许多省份也有，并将他们的住所称为隐修院。由此可知，最初信基督者的教会，正是如今隐修士所愿效法的，即其中无一人有任何私产，无一人富有，无一人贫穷，遗产分给穷人，他们有暇祈祷和咏唱圣咏，还从事教义和苦修实践，事实上正如路加所载，最初在耶路撒冷的信徒们那样。据说在卡里古拉时期，他冒险前往罗马，作为本国使节被派去；第二次到克劳狄那里时，他在同一城市与宗徒伯多禄交谈，并享有他的友谊，因此他也以赞扬美化了马尔谷（伯多禄在亚历山大里亚的门徒）的追随者。此人有许多卓越且难以计数的著作：\WorkTitle{论梅瑟五书}一卷，\WorkTitle{论语言混乱}一卷，\WorkTitle{论自然与发明}一卷，\WorkTitle{论我们感官所欲与所恶之事}一卷，\WorkTitle{论学习}一卷，\WorkTitle{论神圣事物的继承者}一卷，\WorkTitle{论相等与相反之分}一卷，\WorkTitle{论三超德}一卷，\WorkTitle{论为何圣经中许多人的名字被改变}一卷，\WorkTitle{论盟约}两卷，\WorkTitle{论智者的生活}一卷，\WorkTitle{论巨人}一卷，\WorkTitle{梦由天主所遣}五卷，\WorkTitle{出谷纪问答}五卷，\WorkTitle{论会幕与十诫}四卷，以及\WorkTitle{论祭品与许诺或诅咒}、\WorkTitle{论上智安排}、\WorkTitle{论犹太人}、\WorkTitle{论生活方式}、\WorkTitle{论亚历山大}、\WorkTitle{哑兽有理性}、\WorkTitle{每个愚人应为奴隶}，以及我们前面提到的\WorkTitle{论基督徒的生活}，即宗徒式人物的生活，他将其题为\WorkTitle{论度神圣生活者}，因为他们确实默观神圣之事并不断向天主祈祷，此外还有其他类别：\WorkTitle{论农业}两卷，\WorkTitle{论醉酒}两卷。尚有其天才的其他作品未传到我们手中。关于他，希腊人中有一句谚语：\Quote{要么柏拉图效法斐洛，要么斐洛效法柏拉图。}即要么柏拉图追随斐洛，要么斐洛追随柏拉图，他们的思想和语言如此相似。\par

% structure: p0029 source=h2 target=section reason=relative-heading-rank; base=h2

\section{第十二章}

科尔杜巴的\Person{路齐乌斯·安内乌斯·塞内加}，斯多葛派\Person{索提翁}的门徒，诗人\Person{路坎}的舅父，是一位生活极为克己的人。我本不应将他列入圣人之列，若非那些被许多人阅读的\WorkTitle{保禄致塞内加书}和\WorkTitle{塞内加致保禄书}促使我这样做。在这些书信中，当他身为\Person{尼禄}的导师且是当时最有权势之人时，他说他愿意在自己的同胞中拥有\Person{保禄}在基督徒中所处的地位。他在\Person{伯多禄}和\Person{保禄}获得殉道冠冕前两年被\Person{尼禄}处死。\par

% structure: p0031 source=h2 target=section reason=relative-heading-rank; base=h2

\section{第十三章}

\Person{若瑟夫}，\Person{玛提雅}之子，耶路撒冷的司祭，被\Person{韦斯巴芗}及其子\Person{提托}俘虏并流放。他来到罗马后，向父子两位皇帝呈献了七卷\WorkTitle{论犹太人的被掳}，这些书被存于公共图书馆，他因其天才在罗马获得立像的荣誉。他还写了二十卷\WorkTitle{古史}，从创世直到\Person{多米提安}皇帝第十四年，以及两卷\WorkTitle{驳阿皮翁的古史}（\Person{阿皮翁}是亚历山大里亚的语法家，在\Person{卡里古拉}时被外邦人派为使节反对\Person{斐洛}），还写了一部辱骂犹太民族的书。他的另一部书题为\WorkTitle{论一切统治智慧}，其中记载了玛加伯殉道者之死，备受推崇。在他的\WorkTitle{古史}第八卷中，他极其公开地承认基督因奇迹的伟大而被法利塞人所杀，洗者\Person{若翰}确实是先知，耶路撒冷因宗徒\Person{雅各伯}的被杀而被毁。他还以如下方式论及主：\Quote{在同一时期，有一位智者\Person{耶稣}，若可合法地称他为人。因为他是奇妙奇迹的行者，又是那些自由接受真理之人的导师。他也有极多的追随者，既有犹太人也有外邦人，并被相信是基督。当我们首领们的嫉妒导致\Person{比拉多}把他钉在十字架上时，然而起初爱他的人仍坚持到底，因为他第三天活活地向他们显现。许多事，包括这些以及其他奇妙之事，都记载于预言他的先知之歌中，并且从他得名的基督徒教派，至今存在。}\par

% structure: p0033 source=h2 target=section reason=relative-heading-rank; base=h2

\section{第十四章}

加利肋亚省提庇里亚的\Person{犹斯都}也曾试图撰写一部\WorkTitle{犹太事务史}和若干简要的\WorkTitle{圣经注释}，但\Person{若瑟夫}证明他虚妄。据知他与\Person{若瑟夫}同时写作。\par

% structure: p0035 source=h2 target=section reason=relative-heading-rank; base=h2

\section{第十五章}

克肋孟（Clement），宗徒保禄在致斐理伯人书信中曾提及\BibleQuote{与克肋孟以及其他同工，他们的名字都写在生命册上}，他在伯多禄之后为罗马的第四位主教（如果第二位是理诺（Linus），第三位是安纳克莱托（Anacletus）的话），尽管大多数拉丁人认为克肋孟是宗徒之后的第二位。他以罗马教会的名义写了一封特别有价值的\WorkTitle{致格林多教会书}，这封书信在某些地方公开诵读，在我看来，其风格与托名保禄的\WorkTitle{希伯来书}相似，但这封书信与同一封书信不仅在诸多思想上不同，而且在措辞顺序上也有差异，两者在任一方面的相似度都不太大。另有一封托名他的第二封\WorkTitle{书信}，被早期作家拒绝，还有一部长篇\WorkTitle{伯多禄与阿庇翁（Appion）辩论录}，欧瑟比在他的\WorkTitle{教会史}第三卷中拒绝承认其真实性。他在图拉真（Trajan）第三年去世，罗马建造的一座教堂至今仍保存着他名字的记忆。\par

% structure: p0037 source=h2 target=section reason=relative-heading-rank; base=h2

\section{第十六章}

依纳爵（Ignatius），安提约基雅（Antioch）教会的第三位主教（继宗徒伯多禄之后），在图拉真迫害期间被判处喂野兽，被捆绑送往罗马。当他航行至士每拿（Smyrna）时，那里有若望的门徒波利卡普（Polycarp）任主教，他写了一封\WorkTitle{致厄弗所人书}，另一封\WorkTitle{致玛格内西亚人书}，第三封\WorkTitle{致特拉里亚人书}，第四封\WorkTitle{致罗马人书}。离开那里后，他写了\WorkTitle{致非拉德里非亚人书}和\WorkTitle{致士每拿人书}，并特别写了一封\WorkTitle{致波利卡普书}，将安提约基雅的教会托付给他。在这最后一封信中，他为那部我最近翻译的福音书作证，论及基督的位格时说：\Quote{我确实在复活后以肉身看见了他，我相信他就是；}当他来到伯多禄和那些与伯多禄在一起的人面前时，对他们说：\InnerQuote{看！触摸我，看看我，我不是一个无形的神灵}，他们立刻触摸了他并相信了。\newline{}此外，既然我们提到了这样一个人以及他写的\WorkTitle{致罗马人书}，似乎值得提供几段\Quote{摘录}：\Quote{从叙利亚直到罗马，我日夜与野兽搏斗，在陆地与海洋，被绑在十只豹子中间，即看守我的士兵，他们越是受善待就越凶恶。他们的恶行却是我的老师，但我并不因此成义。但愿我享受为我预备的野兽；我祈求它们立刻出现；我甚至会引诱它们快速吞噬我，以免它们像对待一些它们因恐惧而拒绝触碰的人那样对待我。如果它们不愿意，我会强迫它们吞噬我。宽恕我吧，我的孩子们，我知道什么对我有益。现在我才开始做门徒，不渴望任何可见的事物，好让我能获得耶稣基督。让火与十字架、野兽的袭击、骨头的折断、肢体的分割、全身的粉碎、魔鬼的折磨——让这一切都临到我身上，只要我能获得在基督内的喜乐。}\par

当他被判处喂野兽，并因对殉道的热忱而听到狮子吼叫时，他说：\Quote{我是基督的麦粒。我被野兽的牙齿磨碎，好能被找到作为世界之粮。}他在图拉真第十一年被处死，尸骨安葬在安提约基雅城外达夫尼特（Daphnitic）门的墓地里。\par

% structure: p0040 source=h2 target=section reason=relative-heading-rank; base=h2

\section{第十七章}

波利卡普（Polycarp），宗徒若望的门徒并由他祝圣为士每拿（Smyrna）的主教，是亚细亚全区的首领。在那里他见过一些宗徒并曾以他们为师，也见过那些见过主的人。由于有关逾越节日期的一些问题，他在安东尼·庇护（Antoninus Pius）皇帝时代前往罗马，当时阿尼塞托（Anicetus）治理该城的教会。在那里，他将许多因马西昂（Marcion）和瓦伦廷（Valentinus）的劝说而受欺骗的信徒引回信德。当马西昂偶然遇见他并说\Quote{你认识我们吗？}时，他回答说：\Quote{我认识魔鬼的长子。}后来，在马可·安东尼（Marcus Antoninus）和路济乌斯·奥雷利乌斯·科莫杜斯（Lucius Aurelius Commodus）统治期间，在尼禄之后的第四次迫害中，在总督于士每拿开庭、全体民众在竞技场中对他呼喊的情况下，他被焚烧。他写了一封非常有价值的\WorkTitle{致斐理伯人书}，至今仍在亚细亚的聚会中被诵读。\par

% structure: p0042 source=h2 target=section reason=relative-heading-rank; base=h2

\section{第十八章}

帕皮亚（Papias），若望的门徒，亚细亚的希拉波立（Hierapolis）主教，只写了五卷书，题为\WorkTitle{吾主言论集}。他在序言中声称自己不追随各种意见，而是以宗徒为权威，他说：\Quote{我考虑了安德肋和伯多禄所说的，斐理伯、多默、雅各伯、若望、玛窦或我主门徒中的任何其他人所说的，也考虑了主的门徒阿里斯蒂翁（Aristion）和长老若望所说的，与其说我有他们的书可读，不如说直到今天仍能从作者本人那里听到他们的活生生的声音。}从这份名单可以看出，列在门徒中的若望，与他在列举中放在阿里斯蒂翁之后的长老若望不是同一个人。我们之所以这样说，是因为前文提到的意见，即许多人宣称若望最后两封书信不是宗徒的作品，而是长老的作品。\par

据说他出版了一部\WorkTitle{吾主的第二次来临或千年帝国}。依肋内（Irenaeus）、阿波利纳里（Apollinaris）以及其他认为主复活后将以肉身与圣徒一同掌权的人追随他的观点。戴尔都良（Tertullian）在他的\WorkTitle{论信者的希望}一书中，以及佩陶的维克托里诺（Victorinus of Petau）和拉克坦提（Lactantius）也遵循这一观点。\par

% structure: p0045 source=h2 target=section reason=relative-heading-rank; base=h2

\section{第十九章}

\Person{Quadratus}是宗徒的门徒，在雅典的主教\Person{Publius}因对基督的信德而殉道后，他接替了\Person{Publius}的位置，并凭借其信德和勤奋，将因极度恐惧而分散的教会重新聚集起来。当\Person{Hadrian}在雅典过冬以观看厄琉息斯秘仪，并几乎被引入希腊所有的神圣奥秘时，那些仇恨基督徒的人便趁机在未受皇帝指示的情况下迫害信徒。此时，他向\Person{Hadrian}呈献了一部为我们的宗教辩护的著作，这部著作必不可少，充满合理的论证和信德，堪当宗徒的教诲。其中，为表明他所处时代的古老，他说他曾亲眼看见许多被各种疾病折磨的人在犹太地被主治愈，甚至还有一些死人复活。\par

% structure: p0047 source=h2 target=section reason=relative-heading-rank; base=h2

\section{第二十章}

\Person{Aristides}是一位极富辩才的雅典哲学家，也是基督的门徒，但仍穿着哲学家的外衣，他在\Person{Quadratus}呈献著作的同时，也向\Person{Hadrian}呈献了一部著作。这部著作系统地陈述了我们的教义，即一部为基督徒辩护的\WorkTitle{护教篇}，至今仍然存世，并被语文学家视为他天才的丰碑。\par

% structure: p0049 source=h2 target=section reason=relative-heading-rank; base=h2

\section{第二十一章}

\Person{Agrippa}，别号\Person{Castor}，是一位学识渊博的人。他撰写了强有力的驳斥文章，针对异端者\Person{Basilides}所写的二十四卷反对福音的著作，揭露了\Person{Basilides}的所有奥秘，列举了先知\Person{Barcabbas}和\Person{Barchob}，以及其他所有令人听闻生畏的野蛮名字，还有他最高的神\Person{Abraxas}，其名字据希腊人的计算包含了年份。\Person{Basilides}死于\Person{Hadrian}统治时期的\Place{Alexandria}，诺斯替派由此产生。在这动荡的时期，犹太派的首领\Person{Cochebas}用各种酷刑处死了基督徒。\par

% structure: p0051 source=h2 target=section reason=relative-heading-rank; base=h2

\section{第二十二章}

\Person{Hegesippus}生活的年代离宗徒时代不远。他撰写了一部\WorkTitle{历史}，涵盖从我们主的受难到他本人时期的所有教会事件，并收集了许多对读者有益的内容，共五卷，风格质朴，力求模仿他所记述人物的言谈方式。他说他于\Person{Anicetus}（\Person{Peter}之后的第十位主教）时期前往\Place{Rome}，并一直留在那里，直到同城主教\Person{Eleutherius}时期，后者曾是\Person{Anicetus}手下的执事。此外，他驳斥偶像，撰写了一部历史，指出偶像最初是从何种错误中产生的，这部著作也表明了他所活跃的时代。他说：\Quote{他们为死者建造纪念碑和庙宇，正如我们今天所见的，例如为\Person{Hadrian}皇帝的仆人\Person{Antinous}所建的；为了纪念他，还举行了竞赛，建立了一座以他命名的城市，并设立了神殿和司祭。}据说\Person{Hadrian}皇帝曾迷恋\Person{Antinous}。\par

% structure: p0053 source=h2 target=section reason=relative-heading-rank; base=h2

\section{第二十三章}

\Person{Justin}是一位哲学家，穿着哲学家的外衣，是\Place{Palestine}城市\Place{Neapolis}的公民，\Person{Priscus Bacchius}的儿子。他为基督的宗教辛勤工作，以至于他向\Person{Antoninus Pius}及其儿子们和元老院呈献了一部题为\WorkTitle{反异教}的著作，并且不回避十字架的耻辱。他还向这位\Person{Antoninus Pius}的继任者——\Person{Marcus Antoninus Verus}和\Person{Lucius Aurelius Commodus}——呈献了另一部书。他的另一卷\WorkTitle{反异教}也存世，其中他讨论了魔鬼的本性；还有第四部反异教的著作，他命名为\WorkTitle{驳斥}；还有一部\WorkTitle{论天主的至上主权}；另一部他命名为\WorkTitle{诗篇}；还有\WorkTitle{论灵魂}；\WorkTitle{反犹太人对话录}，这是他与犹太人首领\Person{Trypho}的对话；还有著名的\WorkTitle{反Marcion}数卷，\Person{Irenæus}也在其\WorkTitle{反异端}第四卷中提及；还有一部\WorkTitle{反一切异端}，他在致\Person{Antoninus Pius}的\WorkTitle{护教篇}中提及。当他在\Place{Rome}城进行\OriginalTerm{讲学}{\GreekText{διατριβάς}}，并揭发犬儒派\Person{Crescens}——他曾说许多亵渎基督徒的话——贪食和怕死，证明他沉溺于奢侈和情欲之后，最终因\Person{Crescens}的努力和诡计而被指控为基督徒，他为基督流了血。\par

% structure: p0055 source=h2 target=section reason=relative-heading-rank; base=h2

\section{第二十四章}

\Place{Asia}的\Place{Sardis}主教\Person{Melito}向皇帝\Person{Marcus Antoninus Verus}——修辞学家\Person{Fronto}的门徒——呈献了一部为基督教义辩护的书。他还写了其他著作，其中包括以下：\WorkTitle{论逾越节}两卷；\WorkTitle{论先知的生活}一卷；\WorkTitle{论教会}一卷；\WorkTitle{论主日}一卷；\WorkTitle{论信德}一卷；\WorkTitle{论圣咏（？）}一卷；\WorkTitle{论感官}一卷；\WorkTitle{论灵魂与身体}一卷；\WorkTitle{论圣洗}一卷；\WorkTitle{论真理}一卷；\WorkTitle{论基督的诞生}一卷；\WorkTitle{论祂的预言}一卷；\WorkTitle{论好客}一卷；还有一部被称为\WorkTitle{钥匙}的——\WorkTitle{论魔鬼}一卷；\WorkTitle{论若望默示录}一卷；\WorkTitle{论天主的形体性}一卷；以及六卷\WorkTitle{选集}。关于他卓越的演说天才，\Person{Tertullian}在他为支持\Person{Montanus}而反对教会的七卷书中讽刺地说，他被我们许多人视为先知。\par

% structure: p0057 source=h2 target=section reason=relative-heading-rank; base=h2

\section{第二十五章}

\Place{Antioch}教会第六位主教\Person{Theophilus}在皇帝\Person{Marcus Antoninus Verus}统治期间撰写了一部\WorkTitle{反Marcion}，至今存世；还有三卷\WorkTitle{致Autolycus}、一卷\WorkTitle{反Hermogenes的异端}，以及其他简短优雅的论著，非常适合教会的建立。我曾读过他名下的一些注释\WorkTitle{论福音}和\WorkTitle{论撒罗满箴言}，但在我看来，这些作品在风格和语言上与上述著作的优雅和表现力不相符。\par

% structure: p0059 source=h2 target=section reason=relative-heading-rank; base=h2

\section{第二十六章}

\Place{Asia}的\Place{Hierapolis}主教\Person{Apollinaris}活跃于\Person{Marcus Antoninus Verus}统治时期，他向皇帝呈献了一部为基督徒的信德辩护的著名著作。他的另外五卷\WorkTitle{反异教}、两卷\WorkTitle{论真理}和\WorkTitle{反Cataphrygians}亦存世，这些是在\Person{Montanus}与\Person{Prisca}及\Person{Maximilla}刚开始活动时写成的。\par

% structure: p0061 source=h2 target=section reason=relative-heading-rank; base=h2

\section{第二十七章}

科林斯教会的主教狄奥尼修斯，其文才与勤勉如此卓越，以致他不仅以书信教导本城及本省的人民，也教导其他省与城的人民。在这些书信中，有\WorkTitle{致拉栖第梦人书}、\WorkTitle{致雅典人书}、\WorkTitle{致尼科美底亚人书}、\WorkTitle{致克里特人书}、\WorkTitle{致阿马斯特里教会与本都其他教会书}、\WorkTitle{致格诺索斯人及该城主教比尼托书}、\WorkTitle{致罗马人书}（致其主教索特尔）以及\WorkTitle{致克里索弗拉书}（一位圣洁的妇女）。他活跃于马尔库斯·安东尼努斯·维鲁斯与路奇乌斯·奥雷利乌斯·科莫杜斯在位时期。\newline{}\par

% structure: p0063 source=h2 target=section reason=relative-heading-rank; base=h2

\section{第二十八章}

克里特的比尼托，格诺索斯城的主教，致信科林斯的主教狄奥尼修斯，写了一封极为优雅的书信，在其中他教导说，人民不可永远被喂以奶，以免他们尚为婴儿时就被末日追上，而应也喂以固体食物，好使他们迈向属灵的老年。他活跃于马尔库斯·安东尼努斯·维鲁斯与路奇乌斯·奥雷利乌斯·科莫杜斯在位时期。\newline{}\par

% structure: p0065 source=h2 target=section reason=relative-heading-rank; base=h2

\section{第二十九章}

塔提安，在教授修辞学时，于修辞艺术中赢得不少荣耀，是殉道者犹斯定的追随者，只要他不离开其师，便卓然出众。但后来，因辩才膨胀而自大，他创立了一个称为\Quote{戒行派}的新异端。塞维鲁后来增强了这个异端，以致这一派的异端者至今仍被称为塞维鲁派。塔提安除了无数卷著作外，尚存有一部极为成功的著作\WorkTitle{反异教论}，这被认为是他所有作品中最重要的一部。他活跃于马尔库斯·安东尼努斯·维鲁斯与路奇乌斯·奥雷利乌斯·科莫杜斯在位时期。\newline{}\par

% structure: p0067 source=h2 target=section reason=relative-heading-rank; base=h2

\section{第三十章}

克里特的斐理伯，即戈尔提纳城的主教，狄奥尼修斯在致该城教会的书信中提及他。他出版了一部卓越的著作\WorkTitle{反马尔基昂论}，活跃于马尔库斯·安东尼努斯·维鲁斯与路奇乌斯·奥雷利乌斯·科莫杜斯时代。\newline{}\par

% structure: p0069 source=h2 target=section reason=relative-heading-rank; base=h2

\section{第三十一章}

穆萨诺斯，在那些论及教会教义的作者中并非微不足道，于马尔库斯·安东尼努斯·维鲁斯在位时期，写了一本书给某些已脱离教会、转向戒行派异端的弟兄。\newline{}\par

% structure: p0071 source=h2 target=section reason=relative-heading-rank; base=h2

\section{第三十二章}

莫德斯托也在马尔库斯·安东尼努斯与路奇乌斯·奥雷利乌斯·科莫杜斯在位时期，写了一部\WorkTitle{反马尔基昂论}，至今尚存。他的名下还传有其他一些作品，但学者们认为那些是伪作。\newline{}\par

% structure: p0073 source=h2 target=section reason=relative-heading-rank; base=h2

\section{第三十三章}

美索不达米亚的巴尔德撒纳斯被列入杰出人物之中。他最初是瓦伦廷的追随者，后来成为其对手，并亲自创立了一个新异端。他在叙利亚人中享有才华横溢、辩论激烈的声誉。他写了大量反对几乎他那个时代所有异端的著作。其中一部最卓越、最有力的作品是他致马尔库斯·安东尼努斯的\WorkTitle{论命运}，以及许多其他关于\WorkTitle{论迫害}的卷册，他的追随者将这些著作从叙利亚文译成了希腊文。如果译文都显示出如此的力量与光辉，那么原文必定何其伟大。\newline{}\par

% structure: p0075 source=h2 target=section reason=relative-heading-rank; base=h2

\section{第三十四章}

维克多，罗马的第十三位主教，写了\WorkTitle{论逾越节争端}及其他一些小作品。他在塞维鲁皇帝在位期间治理教会十年。\newline{}\par

% structure: p0077 source=h2 target=section reason=relative-heading-rank; base=h2

\section{第三十五章}

依勒内，在高卢里昂教会的主教波提诺属下任长老，因某些教会事务，被该地的殉道者派往罗马作为使节，向主教埃莱乌特里乌斯递交了以其本人名义写的几封值得尊敬的书信。后来波提诺年近九十，为基督接受了殉道的冠冕，依勒内便接替了他的位置。同样可以确定的是，他是上面提到过的司铎兼殉道者波利卡普的门徒。他写了五卷\WorkTitle{反异端论}、一卷短篇\WorkTitle{反异教论}、另一卷\WorkTitle{论纪律}、一封致其弟兄马尔西安的\WorkTitle{论讲道}、一部\WorkTitle{杂论集}；还有致布拉斯特的\WorkTitle{论分裂}、致弗洛里努斯的\WorkTitle{论独一主权}或\WorkTitle{论天主非恶之根源}，以及一部卓越的\WorkTitle{八元论释义}。在这部作品的结尾，他表明自己接近宗徒时代，写道：\Quote{我指着我们的主耶稣基督，并他荣耀的来临——那时他要审判活人死人——恳求你，凡转录此书者，在转录之后，要仔细对照并根据你所转录的范本加以修正；并且你也要同样转录这道誓言，正如你在范本中所找到的那样。}他的其他作品也在流传，例如致罗马主教维克多的\WorkTitle{论逾越节争端}，他在其中警告维克多不要轻易破坏弟兄的合一，如果维克多真的认为亚细亚和东方的许多主教与犹太人一起在新月第十四天庆祝逾越节是该受谴责的话。但即使那些与他们意见不同的人也不支持维克多的意见。他主要活跃于科莫杜斯皇帝在位时期，科莫杜斯接替了马尔库斯·安东尼努斯·维鲁斯执政。\newline{}\par

% structure: p0079 source=h2 target=section reason=relative-heading-rank; base=h2

\section{第三十六章}

庞泰诺，斯多葛学派的哲学家，按照亚历山大里亚的古老风俗——从福音作者马尔谷以来，教会人士一直担任教师——他拥有如此卓越的智慧与学识，无论是在圣经还是在世俗文学方面，以致应那国使节的请求，他被亚历山大里亚主教德默特琉派往印度。在那里，他发现十二宗徒之一的巴尔多禄茂曾按照\BibleBook{玛窦}{福音}宣讲了主耶稣的来临。返回亚历山大里亚时，他将用希伯来文字写成的这部福音带了回来。他的许多关于圣经的注释确实尚存，但他在教会中的活生生的教导更为有益。他在塞维鲁皇帝与绰号为卡拉卡拉的安东尼努斯在位期间教导。\newline{}\par

% structure: p0081 source=h2 target=section reason=relative-heading-rank; base=h2

\section{第三十七章}

罗德，亚细亚人，在罗马随从我们上文提到的塔提安学习圣经，著作颇多，尤其是\WorkTitle{驳马西昂}一书，其中叙述了马西昂派彼此之间以及与教会之间的分歧，并说年迈的异端者阿佩莱斯曾与他辩论，罗德嘲笑阿佩莱斯，因为后者宣称他不知道他所敬拜的天主。在同一部写给卡利斯提翁的书中，他提到自己曾在罗马作塔提安的门徒。他还撰写了优雅的\WorkTitle{论创世六日}以及著名的\WorkTitle{驳弗吕家人}。他活跃于科莫杜斯和塞维鲁在位时期。\par

% structure: p0083 source=h2 target=section reason=relative-heading-rank; base=h2

\section{第三十八章}

克莱孟，亚历山大里亚教会的司铎，是上文提到的潘代努斯的门徒，在其师去世后领导亚历山大里亚的神学院，并担任教理导师。他著作了多部充满雄辩与学问的著名卷册，涉及圣经与世俗文学：其中有\WorkTitle{杂论}八卷、\WorkTitle{纲要}八卷、\WorkTitle{驳异教徒}一卷、\WorkTitle{论教育}三卷、\WorkTitle{论逾越节}、\WorkTitle{论斋戒}以及另一部题为\WorkTitle{富人得救吗？}的书。还有\WorkTitle{论诽谤}一卷、\WorkTitle{论教会法规并驳斥随从犹太人谬误者}一卷，是写给耶路撒冷主教亚历山大的。他在\WorkTitle{杂论}卷中提到了我们上文所述的塔提安著作\WorkTitle{驳异教徒}，以及一位名叫卡西安的作者的\WorkTitle{编年史}，这是我未能找到的一部作品。他还提到某些犹太作家驳斥异教徒的作品，如阿里斯托布鲁斯、德米特里和犹波莱摩，他们效仿约瑟夫斯，主张梅瑟和犹太民族的优先地位。有一封耶路撒冷主教亚历山大的信，他后来与纳尔基索共同治理该教会，关于宣信者阿斯克勒皮亚德斯的任命，致安提约基雅人表示祝贺，在信末他说：\Quote{尊敬的弟兄们，我已将这些著作托付给有福的司铎克莱孟带给你们，他是一位杰出而受称许的人，你们也认识他，如今你们将更为了解他；他因天主特别的眷顾来到这里，巩固并扩展了天主的教会。}奥利振是他的门徒。他活跃于塞维鲁及其子安托尼努斯在位时期。\par

% structure: p0085 source=h2 target=section reason=relative-heading-rank; base=h2

\section{第三十九章}

密提亚德，罗德在反对孟他努、普里斯卡和马克西米拉的著作中曾提到他，他写了大量反对这些人的作品，以及其他\WorkTitle{驳异教徒与犹太人}的书籍，并向当时在位的皇帝呈递了一份\WorkTitle{护教书}。他活跃于马尔库斯·安托尼努斯和科莫杜斯在位时期。\par

% structure: p0087 source=h2 target=section reason=relative-heading-rank; base=h2

\section{第四十章}

阿波罗尼乌斯，一位极具才华的人，写了一部反对孟他努、普里斯卡和马克西米拉的著名长篇著作，其中断言孟他努及其疯狂的女先知是上吊而死，并记载了许多其他事情，其中包括以下关于普里斯卡和马克西米拉的话：\Quote{若她们否认接受了礼物，就让她们承认接受礼物的人不是先知，我将用上千证人证明她们接受了礼物，因为先知是凭其他果实才显明确实是先知的。告诉我，先知会染头发吗？先知会用锑粉涂眼睑吗？先知会穿着华服和宝石吗？先知会掷骰子和玩棋盘吗？他会收取利息吗？让她们回答这是否应当允许，我的任务就是证明她们做了这些事。}他在同一本书中说，他写这部作品时，正是卡塔弗里吉亚异端开始后的第四十年。德尔图良在他反对教会所写的六卷\WorkTitle{论神魂超拔}之外，又加上了第七卷，特别针对阿波罗尼乌斯，企图为阿波罗尼乌斯所驳斥的一切辩护。阿波罗尼乌斯活跃于科莫杜斯和塞维鲁在位时期。\par

% structure: p0089 source=h2 target=section reason=relative-heading-rank; base=h2

\section{第四十一章}

塞拉皮翁，在科莫杜斯皇帝在位第十一年被祝圣为安提约基雅主教，他写了一封信给卡里库斯和彭提乌斯，论及孟他努异端，其中说：\Quote{此外，为使你们知道这虚假教义——即新预言教义——的疯狂被全世界所遣责，我已将亚细亚希拉波利斯最神圣的主教阿波利纳里的信函寄给你们。}他还写了一部著作给多米努斯，后者在迫害时期投靠了犹太人；另一部著作论及以伯多禄之名流传的福音，致基里基亚的罗森斯教会，该教会因阅读此书而陷入异端。此外，他还有多处短函，与其作者之刻苦生活相称。\par

% structure: p0091 source=h2 target=section reason=relative-heading-rank; base=h2

\section{第四十二章}

阿波罗尼乌斯，科莫杜斯皇帝时期的罗马元老院议员，因一名奴隶告发其为基督徒，获准为自己的信仰辩护，并写了一部卓越的著作，在元老院宣读；然而，按照元老院的决议，他仍因他们当中一条古老的法律被斩首殉道，该法律规定，凡曾受其审判的基督徒，若不悔改，不得释放。\par

% structure: p0093 source=h2 target=section reason=relative-heading-rank; base=h2

\section{第四十三章}

德奥斐洛，巴勒斯坦凯撒勒雅（该城旧称斯特拉托塔）的主教，在塞维鲁皇帝在位期间，与其他主教共同撰写了一封极有用处的教务会议信函，反对那些随同犹太人在当月第十四日庆祝逾越节的人。\par

% structure: p0095 source=h2 target=section reason=relative-heading-rank; base=h2

\section{第四十四章}

巴基鲁斯，格林多的主教，在同一位塞维鲁皇帝在位期间享有盛誉，他代表亚该亚所有主教撰写了一部优美著作\WorkTitle{论逾越节}。\par

% structure: p0097 source=h2 target=section reason=relative-heading-rank; base=h2

\section{第四十五章}

厄弗所主教\Person{波利克拉特斯}，与亚细亚的其他主教，他们按照某种古代习俗与犹太人在该月的第十四日庆祝逾越节，写了一封反对罗马主教\Person{维克托}的会议书信，他在其中说他遵循宗徒\Person{若望}和古代先贤的权威。我们从这封信中摘录以下简短的引文：\Quote{因此，我们按照惯例不可侵犯地庆祝这一天，既不添加什么，也不删减什么，因为在亚细亚安息着那些将在主的日子复活的最伟大圣人们的遗骸——当主带着威严从天降临，并使所有圣人都活过来时——我指的是十二宗徒之一\Person{斐理伯}，他安息在\Place{希拉波利斯}，还有他的两个女儿，她们至死守贞，以及他的另一个女儿，她在\Place{厄弗所}充满圣神而死。还有\Person{若望}，他曾靠在主的胸膛，是主的大司祭，额上戴着金牌，既是殉道者又是圣师，在\Place{厄弗所}安眠；\Person{波利卡普}主教兼殉道者死在\Place{斯米纳}。\Place{欧墨尼亚}的\Person{特拉塞亚斯}，也是主教兼殉道者，安息在同一个\Place{斯米纳}。何须提及\Person{萨加里斯}，主教兼殉道者，他安睡在\Place{劳狄刻雅}，以及可敬的\Person{帕皮鲁斯}和在圣神内的阉者\Person{梅利托}，他永远事奉主，在\Place{撒尔德}安葬，在那里等待基督来临时他的复活。这些人都遵守在该月第十四日庆祝逾越节，丝毫不偏离福音传统，并遵循教会规范。我，\Person{波利克拉特斯}，你们众仆人中最小的一位，按照我亲属的教导（我亦遵循此教导，因我有七位亲属是主教，而我是第八位），也总是在犹太民族庆祝除去旧酵时庆祝逾越节。所以，弟兄们，我在主内已六十五岁，并由来自世界各地的许多弟兄教导，查考了全部圣经，我必不怕那些恐吓我们的人，因为我的前任们说过：\InnerQuote{应当服从天主，而不服从人。}}我引用这些，为通过一个小例子表明此人的才智和权威。他在\Person{塞维鲁}皇帝统治时期兴盛，与\Place{耶路撒冷}的\Person{纳尔奇苏斯}同时代。\par

% structure: p0099 source=h2 target=section reason=relative-heading-rank; base=h2

\section{第四十六章}

\Person{赫拉克利图斯}在\Person{科莫杜斯}和\Person{塞维鲁}统治时期，撰写了关于\WorkTitle{宗徒大事录}和\WorkTitle{书信}的注释。\par

% structure: p0101 source=h2 target=section reason=relative-heading-rank; base=h2

\section{第四十七章}

\Person{马克西穆斯}在同样几位皇帝统治下，在一部卓越的著作中提出了著名的问题：\Emph{恶的起源是什么？}以及\Emph{物质是否由天主创造。}\par

% structure: p0103 source=h2 target=section reason=relative-heading-rank; base=h2

\section{第四十八章}

\Person{坎迪杜斯}在以上提及的皇帝统治下，出版了极为出色的论著\WorkTitle{论创造六天}。\par

% structure: p0105 source=h2 target=section reason=relative-heading-rank; base=h2

\section{第四十九章}

\Person{阿皮翁}在\Person{塞维鲁}皇帝统治下，也撰写了论著\WorkTitle{论创造六天}。\par

% structure: p0107 source=h2 target=section reason=relative-heading-rank; base=h2

\section{第五十章}

\Person{塞克斯图斯}在\Person{塞维鲁}皇帝统治下，写了一部书\WorkTitle{论复活}。\par

% structure: p0109 source=h2 target=section reason=relative-heading-rank; base=h2

\section{第五十一章}

\Person{阿拉比亚努斯}在同样一位皇帝统治下，出版了一些与基督信仰教义相关的小作品。\par

% structure: p0111 source=h2 target=section reason=relative-heading-rank; base=h2

\section{第五十二章}

\Person{犹达斯}详细讨论了\Person{达尼尔}书中提到的七十周，并写了一部前代历史的\WorkTitle{编年史}，他将之续写至\Person{塞维鲁}的第十年。他因这部作品被证明有错误，因为他预言假基督的来临将在他那个时期，但这只是因为迫害的剧烈似乎预示着世界的终结。\par

% structure: p0113 source=h2 target=section reason=relative-heading-rank; base=h2

\section{第五十三章}

\Person{戴尔都良}司铎，如今被视为继\Person{维克托}和\Person{阿波罗尼乌斯}之后拉丁作家中的首要人物，他来自非洲省的\Place{迦太基}城，是总督或百夫长之子，一位性格尖锐而有力的人。他主要兴盛于\Person{塞维鲁}皇帝和\Person{安东尼努斯·卡拉卡拉}皇帝统治时期，写了许多卷著作，我们略过不提，因为它们为多数人所熟知。我本人见过一位名叫\Person{保禄}的老者，他是意大利城镇\Place{康科尔迪亚}人，他本人年轻时曾是已年长的可敬的\Person{西彼廉}的秘书。他说他亲眼见过\Person{西彼廉}的习惯：没有一天不读\Person{戴尔都良}，而且他常对他说：\Quote{给我师傅}，意思是指\Person{戴尔都良}。他是教会司铎直至中年，后来因罗马教会圣职人员的嫉妒和辱骂而被赶走，他堕入\Person{蒙塔努斯}的教义，并在自己的许多书中提到新预言。\par

此外，他还直接针对教会编写了以下著作：\WorkTitle{论端庄}、\WorkTitle{论迫害}、\WorkTitle{论斋戒}、\WorkTitle{论一夫一妻}、六卷\WorkTitle{论神往}，以及他写的第七卷\WorkTitle{反阿波罗尼乌斯}。据说他活到衰迈的晚年，并写了许多小作品，但这些已不存世。\par

% structure: p0116 source=h2 target=section reason=relative-heading-rank; base=h2

\section{第五十四章}

奥力振，绰号亚当曼提斯，在塞维鲁·佩蒂纳克斯在位第十年基督徒遭受迫害时，他的父亲莱昂尼达斯为基督接受了殉道者的荣冠，他本人约十七岁，与六个兄弟和寡母陷入贫困，因为他们的财产因承认基督而被没收。年仅十八岁时，他承担了在\Place{亚历山大里亚}分散的教会中教导教理受教者的工作。后来，由该城主教\Person{德默特琉}任命，接替长老\Person{克莱孟}，他兴盛多年。当他已届中年时，由于\Place{阿哈雅}的教会为许多异端所困扰，他持教会信函，取道\Place{巴勒斯坦}前往\Place{雅典}，并由\Place{凯撒勒雅}和\Place{耶路撒冷}的主教\Person{特奥克蒂斯托}和\Person{亚历山大}擢升为长老，因而触怒了\Person{德默特琉}，他狂怒至极，到处写信损害他的声誉。众所周知，在去凯撒勒雅之前，他曾在\Place{罗马}，在主教\Person{则斐琳}治下。他立即返回亚历山大里亚后，使长老\Person{赫拉克拉斯}（他仍穿着哲学家的服饰）成为他在教理学校的助手。赫拉克拉斯在德默特琉之后成为亚历山大里亚教会的主教。奥力振的光荣何等伟大，从以下事实可见：凯撒勒雅主教\Person{费尔米利亚诺}与所有\Place{卡帕多细亚}的主教们邀请他访问，并款待了他很长时间。后来，他去巴勒斯坦朝圣，来到凯撒勒雅，长时间受教于奥力振，学习圣经。也可见于：应亚历山大皇帝的母亲\Person{玛迈亚}（一位虔诚的妇人）的请求，他去\Place{安提约基雅}，在那里受到极大尊敬，并致信皇帝\Person{斐理伯}（他是罗马统治者中第一位成为基督徒的）及其母亲，这些信件至今尚存。谁不知道他如此勤勉于圣经研究，以至于违背其时代和民族的精神，学习了希伯来语，并且以七十贤士译本为基础，收集了其他译本汇成一部作品，即\Person{阿奎拉}、本都人\Person{普罗塞利特}、厄俾尼派\Person{特奥多提安}以及同派\Person{辛马库}的译本，后者还撰写了关于\BibleBook{玛窦}{福音}的注释，试图从中建立自己的教义。此外，他还以极大勤奋寻得了第五、第六和第七种译本（这些译本我们也有从他的图书馆得来），并与其他版本作了比较。由于我已在我写给\Person{保拉}的书信集以及我针对\Person{瓦罗}作品所写的信中列出了他的著作清单，现在暂且不提，然而我不忘记提及他不朽的天才：他精通辩证法、几何学、算术、音乐、语法和修辞学，并教导所有哲学学派，以至于他甚至有勤奋的世俗文学学生，每日为他们讲课，聚集到他那里的人群令人惊叹。他接纳他们，希望藉此世俗文学的工具，将他们在基督信仰上建立起来。\par

无需详述\Person{德西乌}（他对\Person{斐理伯}的宗教疯狂而杀害了他）手下那场针对基督徒的迫害的残酷——在那场迫害中，罗马教会的主教\Person{法比安}在罗马遇难，\Place{耶路撒冷}和\Place{安提约基雅}教会的主教\Person{亚历山大}和\Person{巴比拉斯}因承认基督而被监禁。若有人想知道关于奥力振立场的情况，他首先可以从奥力振自己在那次迫害后写给不同人的书信中清楚了解，其次可从\Place{凯撒勒雅}的\Person{欧瑟比}所著的\WorkTitle{教会史}第六卷以及他为同一奥力振所写的六卷书中了解。\par

他活到\Person{加卢斯}和\Person{沃鲁西安}时代，即他的六十九岁，死于\Place{提洛}，也葬在该城。\par

% structure: p0120 source=h2 target=section reason=relative-heading-rank; base=h2

\section{第五十五章}

\Person{阿摩尼乌}，一位才华横溢、哲学学识丰富的人，同时期在\Place{亚历山大里亚}声名卓著。在他众多杰出的天才作品中，有他精心编撰的\WorkTitle{论梅瑟与耶稣的和谐}以及他编成的\WorkTitle{福音合参}，后来\Place{凯撒勒雅}的\Person{欧瑟比}沿用了它。\Person{波菲利}错误地指控他，说他先作基督徒，后又重新成为异教徒，但可以肯定他终生都保持着基督徒身份。\par

% structure: p0122 source=h2 target=section reason=relative-heading-rank; base=h2

\section{第五十六章}

\Person{安博罗削}，最初是马尔基翁派，后来被奥力振纠正，在教会中任执事，并作为主的宣信者而光荣卓越。奥力振的\WorkTitle{论殉道}一书就是写给他和长老\Person{普罗托克提图}的。奥力振在他的勤奋、资助和坚持下口述了大量著作。安博罗削本人，正如一位高尚天性之人所应有的，并非没有相当的文学才华，他的书信可向奥力振表明。他死于奥力振之前，受到许多人谴责，因为他是富人，临终时没有在遗嘱中记念他年迈贫困的朋友。\par

% structure: p0124 source=h2 target=section reason=relative-heading-rank; base=h2

\section{第五十七章}

\Person{特里弗}，奥力振的门徒，奥力振的一些现存书信是写给他的。他精通圣经，他的许多作品在这里和那里都显示了这一点，尤其是他编著的\WorkTitle{论申命纪中的红母牛}和\WorkTitle{论半份}，即\Person{亚伯拉罕}在\BibleBook{创世}{纪}中献鸽子和斑鸠时所献上的。\par

% structure: p0126 source=h2 target=section reason=relative-heading-rank; base=h2

\section{第五十八章}

\Person{米努齐乌斯·斐理克斯}，\Place{罗马}一位杰出的辩护人，写了一篇对话，代表一位基督徒和一位外邦人之间的讨论，题为\WorkTitle{屋大维}；另有题为\WorkTitle{论命运}或\WorkTitle{反对占星家}的作品以他的名义流传，但这虽然是天才之作，在我看来其风格与上述作品不符。\Person{拉克坦提乌斯}在他的作品中也提到过这位米努齐乌斯。\par

% structure: p0128 source=h2 target=section reason=relative-heading-rank; base=h2

\section{第五十九章}

伽尤斯，罗马的主教，在责斐利诺时代，即在塞维鲁之子安多尼诺统治期间，发表了一篇非常著名的辩论\WorkTitle{反对普罗库鲁斯}，他是蒙塔努斯的追随者，指证他在辩护新预言时的轻率；在同一卷书中，他也只列举了保禄的十三封书信，他说第十四封，即现在称为\WorkTitle{致希伯来人书}的，并非保禄所写，且在罗马至今不被认为是宗徒保禄的作品。\par

% structure: p0130 source=h2 target=section reason=relative-heading-rank; base=h2

\section{第六十章}

贝里鲁斯，阿拉伯地区波斯脱拉的主教，在光荣地治理教会一段时间后，最终陷入否认基督在降生成人之前存在的异端。经奥力振纠正后，他写了各种短篇作品，尤其是书信，在其中他感谢奥力振。奥力振写给他的书信也存世，还有奥力振与贝里鲁斯之间的对话，其中讨论了异端问题。他在玛迈亚之子亚历山大、以及继承其权力的马克西米努斯和戈尔迪安统治期间闻名。\par

% structure: p0132 source=h2 target=section reason=relative-heading-rank; base=h2

\section{第六十一章}

希玻里托斯，某教会的主教（我未能得知该城之名），撰写了\WorkTitle{逾越节的计算与年代列表}，他一直推算到亚历山大皇帝的第一年。他还讨论了十六年周期，希腊人称之为\OriginalTerm{十六年周期}{\GreekText{ἐκκαιδεκαετηρίδα}}，并为欧瑟比提供了线索，后者就同一逾越节编制了一个十九年周期，即\OriginalTerm{十九年周期}{\GreekText{ἐννεακαιδεκαετηρίδα}}。他写了一些圣经注释，其中包括以下著作：\WorkTitle{论创世六日}、\WorkTitle{论出谷纪}、\WorkTitle{论雅歌}、\WorkTitle{论创世纪}、\WorkTitle{论匝加利亚}、\WorkTitle{论圣咏}、\WorkTitle{论依撒意亚}、\WorkTitle{论达尼尔}、\WorkTitle{论默示录}、\WorkTitle{论箴言}、\WorkTitle{论训道篇}、\WorkTitle{论撒乌耳}、\WorkTitle{论比通尼撒}、\WorkTitle{论假基督}、\WorkTitle{论复活}、\WorkTitle{驳马尔基盎}、\WorkTitle{论逾越节}、\WorkTitle{驳一切异端}，以及一篇劝勉\WorkTitle{赞颂我们的主与救主}，在其中他表明自己是在奥力振面前于教会中发言。我们曾提及的安布罗修，他从马尔基盎的异端被奥力振引向真信仰，他激励奥力振为了与希玻里托斯竞争而撰写圣经注释，为他提供了七名甚至更多的秘书及其费用，以及同等数量的抄写员，并且更为惊人的是，他以难以置信的热忱每天要求奥力振工作，为此奥力振在一封书信中称他为\Quote{工头}。\par

% structure: p0134 source=h2 target=section reason=relative-heading-rank; base=h2

\section{第六十二章}

亚历山大，卡帕多细亚的主教，渴望探访圣地，来到耶路撒冷，当时该城的主教纳尔奇苏斯已年迈，治理教会。纳尔奇苏斯和他的许多神职人员得到启示，次日早晨会有一位主教进城，他将成为司祭宝座的助手。事情正如所预言的发生了，巴勒斯坦的所有主教聚集在一起，纳尔奇苏斯本人尤其恳切，亚历山大接掌了耶路撒冷教会的舵柄。在他写给安提诺人论及教会和平的一封书信的末尾，他说：\Quote{纳尔奇苏斯，他此前在这里担任主教，现在与我一同藉他的祈祷履行职责，他大约一百一十六岁，问候你们，并与我一同恳求你们同心合意。}他还写了另一封书信\WorkTitle{致安提约基雅人}，由我们上述的亚历山大里亚司铎克莱孟代笔；另一封\WorkTitle{致奥力振}，以及\WorkTitle{为奥力振辩护驳德默特琉}，这是因德默特琉作证说他已按立奥力振为司铎而引发的。他还有其他致不同人士的书信。在德西乌斯统治下的第七次迫害中，当时安提约基雅的巴比拉被处死，亚历山大被带到凯撒勒雅并囚禁在狱中，他因宣认基督而获得了殉道的荣冠。\par

% structure: p0136 source=h2 target=section reason=relative-heading-rank; base=h2

\section{第六十三章}

犹利乌·阿非利加努斯，其五卷\WorkTitle{编年史}尚存，在继玛克里努斯之后的马尔库斯·奥雷略·安多尼诺统治期间，受命重建厄玛乌城，该城后来被称为尼科波利斯。他有一封致奥力振的书信\WorkTitle{论苏撒纳的问题}，其中争辩说这个故事并未包含在希伯来原文中，且与希伯来语源学关于\Quote{prinos 和 prisai}、\Quote{schinos 和 schisai}的文字游戏不一致。奥力振为此写了一封博学的回信。另有一封信\WorkTitle{致亚里斯提德}存世，其中他详细讨论了我们的救主在玛窦和路加所记载的族谱中出现的差异。\par

% structure: p0138 source=h2 target=section reason=relative-heading-rank; base=h2

\section{第六十四章}

革米努斯，安提约基雅教会的司铎，创作了几部天才的著作，活跃于亚历山大皇帝和其本城主教则本努斯时代，尤其是在赫拉克拉斯被祝圣为亚历山大里亚教会主教之时。\par

% structure: p0140 source=h2 target=section reason=relative-heading-rank; base=h2

\section{第六十五章}

特奥多鲁斯，后来被称为额我略，本都地区新凯撒勒雅的主教，在非常年轻的时候，与他的兄弟阿特诺多鲁斯一起，从卡帕多细亚前往贝里图斯，然后到巴勒斯坦的凯撒勒雅，学习希腊和拉丁文学。奥力振看到这两人非凡的天赋，便鼓励他们学习哲学，在教学中逐渐引入对基督的信仰，并使他们也成为自己的追随者。因此，在他们受教五年后，奥力振打发他们回到母亲身边。特奥多鲁斯在离开时写了一篇向奥力振致谢的颂词，在奥力振本人出席的大会上发表。这篇颂词至今尚存。\par

他还写了一篇简短但极有价值的\WorkTitle{训道篇}释义，流传的报告还提到他其他的书信，但尤其提到他作为主教所行的征兆和奇事，为教会带来了极大的荣耀。\par

% structure: p0143 source=h2 target=section reason=relative-heading-rank; base=h2

\section{第六十六章}

罗马主教科尔乃略，西彼廉有八封信函致他。他致信给安提约基雅教会主教法比乌斯，论及 \WorkTitle{罗马、意大利及阿非利加的会议}，另一封 \WorkTitle{论诺瓦提安及背弃信德者}，第三封 \WorkTitle{论会议法案}，第四封致同一法比乌斯的长信，详述诺瓦提安异端的原因及其绝罚。他在加鲁斯与沃鲁西安执政期间治理教会两年。他为基督接受了殉道的荣冠，由路爵继任。\par

% structure: p0145 source=h2 target=section reason=relative-heading-rank; base=h2

\section{第六十七章}

阿非利加的西彼廉，起初以修辞学教师闻名，后来在长老采齐略（他从其获得姓氏）的劝说下成为基督徒，并将所有财产施与穷人。不久后他被引入长老团，并被立为迦太基主教。无需列举他天才的著作，因为它们比太阳更显明。\par

他在瓦肋黎安与加里恩努斯皇帝治下，第八次迫害中被处死，与科尔乃略在罗马被处死的同一天，但不在同一年。\par

% structure: p0148 source=h2 target=section reason=relative-heading-rank; base=h2

\section{第六十八章}

本狄，西彼廉的执事，与他一同流亡至死，留下了一部著名的著作 \WorkTitle{论西彼廉的生平与死亡}。\par

% structure: p0150 source=h2 target=section reason=relative-heading-rank; base=h2

\section{第六十九章}

亚历山大里亚主教狄约尼削，曾作为长老在赫拉克拉斯手下负责教理学校，是奥力振最杰出的学生。他赞同西彼廉及阿非利加主教会议关于重新为异端者施洗的教义，致信多人，这些信函至今流传。他致信安提约基雅教会主教法比乌斯，作 \WorkTitle{论补赎} 一书；另经由依玻理之手致信 \WorkTitle{致罗马人}；致西斯笃（斯德望的继任者）两封信；致罗马教会长老斐理孟与狄约尼削两封信；另致同一狄约尼削（后为罗马主教）一信；并致诺瓦提安一信，论及他们声称诺瓦提安违背其意愿被祝圣为罗马主教一事。此信的开头如下：\Quote{狄约尼削致其弟兄诺瓦提安，祝安。你若真如所言，是并不情愿被祝圣的，当你自愿退位时，便能证此。}\par

他还有另一封书信 \WorkTitle{致狄约尼削与狄迪默}，以及许多 \WorkTitle{论逾越节的节庆信函}，以演讲风格写成；另有一封致亚历山大里亚教会的 \WorkTitle{论流亡}，一封 \WorkTitle{致埃及主教希埃拉克}，还有 \WorkTitle{论死亡}、\WorkTitle{论安息日}、\WorkTitle{论体育馆}，以及一封 \WorkTitle{致赫尔马孟} 及其他 \WorkTitle{论德西乌斯迫害} 的书信；还有两卷 \WorkTitle{反驳内波主教}，后者在其著作中主张身体内的千年王国。他于此间勤奋地讨论了 \BibleBook{若望}{默示录}，并著有 \WorkTitle{反驳撒贝里}、\WorkTitle{致贝尔尼切主教阿蒙}、\WorkTitle{致特肋斯福禄}、\WorkTitle{致欧弗拉诺}，以及四卷 \WorkTitle{致罗马主教狄约尼削}；致劳狄刻雅人 \WorkTitle{论补赎}；致奥力振 \WorkTitle{论殉道}；致亚美尼亚人 \WorkTitle{论补赎}；另 \WorkTitle{论过犯的次序}；致弟茂德 \WorkTitle{论本性}；致欧弗拉诺 \WorkTitle{论诱惑}；还有许多致巴西里德的信函，其中一封声称他也开始为 \BibleBook{}{训道篇} 撰写注释。他去世前数日写的反对萨莫萨特的保禄的著名书信也在流传。他在加里恩努斯在位第十二年去世。\par

% structure: p0153 source=h2 target=section reason=relative-heading-rank; base=h2

\section{第七十章}

罗马长老诺瓦提安，企图篡夺科尔乃略所占的司祭宝座，并建立了诺瓦提安派（希腊人称其为洁净派）的教义，拒绝接纳悔改的背教者。此教义的创始人诺瓦图，原是西彼廉的长老。他著有 \WorkTitle{论逾越节}、\WorkTitle{论安息日}、\WorkTitle{论割损}、\WorkTitle{论司祭职}、\WorkTitle{论祈祷}、\WorkTitle{论犹太人的食物}、\WorkTitle{论热心}、\WorkTitle{论阿塔鲁}，以及许多其他作品，尤其是巨著 \WorkTitle{论天主圣三}，可视为戴尔都良作品的缩影，许多人误将其归于西彼廉。\par

% structure: p0155 source=h2 target=section reason=relative-heading-rank; base=h2

\section{第七十一章}

安提约基雅教会极具天赋的长老马尔基翁，曾在同城极为成功地教授修辞学，与萨莫萨特的保禄（此人作为安提约基雅教会主教引入了阿尔特蒙的教义）进行了一场辩论，由速记员记录下来。这一对话录至今尚存；另有一封由他以会议名义所写的长篇书信，致罗马主教狄约尼削与亚历山大里亚主教玛克西穆。他在克劳狄与奥勒良努斯在位时期活跃。\par

% structure: p0157 source=h2 target=section reason=relative-heading-rank; base=h2

\section{第七十二章}

美索不达米亚主教阿尔赫劳，用叙利亚语撰写了与自波斯而来的摩尼进行辩论的书卷。此书被译为希腊文，为多人所藏。\par

他在普罗布斯皇帝（继奥勒良努斯与塔西佗之后）在位时期活跃。\par

% structure: p0160 source=h2 target=section reason=relative-heading-rank; base=h2

\section{第七十三章}

叙利亚劳狄刻雅主教、亚历山大里亚的安纳托里，在普罗布斯与卡鲁斯皇帝在位时期活跃，在算术、几何、天文、文法、修辞与辩证方面学问惊人。从他所著的 \WorkTitle{论逾越节} 一书以及十卷 \WorkTitle{算术原理} 中，可窥其才华之伟大。\par

% structure: p0162 source=h2 target=section reason=relative-heading-rank; base=h2

\section{第七十四章}

佩托维主教维多利诺，对拉丁文与希腊文同样不熟悉。因此，他的作品虽思想高贵，文体却较逊色。其作品如下：\WorkTitle{论创世纪}、\WorkTitle{论出谷纪}、\WorkTitle{论肋未纪}、\WorkTitle{论依撒意亚}、\WorkTitle{论厄则克耳}、\WorkTitle{论哈巴谷}、\WorkTitle{论训道篇}、\WorkTitle{论雅歌}、\WorkTitle{论若望默示录}注释、\WorkTitle{反驳一切异端}及许多其他作品。他最终接受了殉道的荣冠。\par

% structure: p0164 source=h2 target=section reason=relative-heading-rank; base=h2

\section{第七十五章}

庞菲禄司铎，凯撒勒雅主教欧瑟比的赞助人，对圣经文学如此热忱，以至于亲手抄写了奥力振的大部分著作，这些著作至今仍保存在凯撒勒雅的图书馆中。我拥有二十五卷奥力振的注释，由他亲手书写，题为\WorkTitle{论十二位先知}，我如此喜悦地珍藏并守护它们，以至于自认为拥有了克罗伊斯般的财富。如果拥有一封殉道者的书信是如此喜悦，那么拥有数千行仿佛用他的鲜血写成的文字，该是何等更大的喜悦！他在欧瑟比之前撰写了\WorkTitle{为奥力振辩护}，并在马克西米努斯迫害期间在巴勒斯坦的凯撒勒雅被处死。\par

% structure: p0166 source=h2 target=section reason=relative-heading-rank; base=h2

\section{第七十六章}

彼厄里乌斯，在卡鲁斯和戴克里先统治时期，亚历山大里亚教会的司铎，当同教会的塞奥纳斯担任主教时，他极其成功地教导民众，达到了如此优雅的文辞，并发表了诸多论及各类主题的论文（至今尚存），因而被称为\Quote{小奥力振}。他以自律著称，甘愿守神贫，精通辩证法。迫害之后，他在罗马度过了余生。现存他的一篇长论\WorkTitle{论欧瑟亚先知}，从内容看，似乎是在逾越节前夕宣讲的。\par

% structure: p0168 source=h2 target=section reason=relative-heading-rank; base=h2

\section{第七十七章}

路济安，极具才华之人，安提约基雅教会的司铎，对圣经研究如此勤奋，以至于至今某些圣经抄本仍冠以路济安之名。他的著作\WorkTitle{论信德}以及致多人的短篇\WorkTitle{书信}尚存。他在马克西米努斯迫害期间，因宣认基督而在尼科美底亚被处死，埋葬于俾提尼亚的海伦诺波利斯。\par

% structure: p0170 source=h2 target=section reason=relative-heading-rank; base=h2

\section{第七十八章}

菲肋亚斯，是那座被称为特穆伊斯的埃及城市的居民，出身贵族，家产丰厚，成为主教后，撰写了一篇精美的作品，赞美殉道者，并驳斥那位试图强迫他献祭的审判官，在同一次迫害中（路济安在尼科美底亚被处死的同一时期），他为基督而被斩首。\par

% structure: p0172 source=h2 target=section reason=relative-heading-rank; base=h2

\section{第七十九章}

阿尔诺比乌斯，在戴克里先统治时期，于非洲的西卡城是一位非常成功的修辞学教师，并撰写了多卷\WorkTitle{驳异民}的著作，这些著作随处可见。\par

% structure: p0174 source=h2 target=section reason=relative-heading-rank; base=h2

\section{第八十章}

菲尔米安，也称为拉克坦提乌斯，是阿尔诺比乌斯的学生。在戴克里先统治时期，他与语法学家弗拉维乌斯一同被召至尼科美底亚（弗拉维乌斯的诗作\WorkTitle{论医学}尚存），并在那里教授修辞学。由于学生稀少（因为这是一座希腊城市），他转而从事写作。我们拥有他年轻时在非洲所写的\WorkTitle{宴会}，以及一篇从非洲至尼科美底亚的\WorkTitle{旅程记}（用六音步诗行写成），另一本称为\WorkTitle{语法学家}的书，还有一本极其优美的\WorkTitle{论天主的义怒}，以及\WorkTitle{神圣原理：驳异民}七卷，同一著作的\WorkTitle{纲要}一卷（无标题），另有\WorkTitle{致阿斯克勒庇亚得斯}两卷、\WorkTitle{论迫害}一卷、\WorkTitle{致普罗布斯书信}四卷、\WorkTitle{致塞维鲁书信}两卷、\WorkTitle{致其弟子德默特琉书信}两卷，以及致同一人的\WorkTitle{论天主的工程或人的受造}一卷。他在极其年迈时，担任了高卢的克里斯普斯·凯撒（君士坦丁之子，后来被其父处死）的导师。\par

% structure: p0176 source=h2 target=section reason=relative-heading-rank; base=h2

\section{第八十一章}

欧瑟比，巴勒斯坦凯撒勒雅的主教，勤于研究圣经，并与殉道者庞菲禄一同成为圣经的最勤勉研究者。他出版了众多著作，其中包括：\WorkTitle{福音论证}二十卷、\WorkTitle{福音预备}十五卷、\WorkTitle{论主显}五卷、\WorkTitle{教会史}十卷、\WorkTitle{世界史编年}及该书的\WorkTitle{纲要}。另有\WorkTitle{论福音书的歧异}十卷、\WorkTitle{论依撒意亚}十卷，以及\WorkTitle{驳波斐利}（此人当时在西西里写作，如某些人所认为）二十五卷，还有\WorkTitle{论题}一卷、\WorkTitle{为奥力振辩护}六卷、\WorkTitle{论庞菲禄生平}三卷、其他短篇\WorkTitle{论殉道者}、极其博学的\WorkTitle{一百五十篇圣咏注释}等。他主要活跃于君士坦丁大帝和君士坦提乌斯统治时期。他因与殉道者庞菲禄的友谊而被称为\Quote{庞菲禄}。\par

% structure: p0178 source=h2 target=section reason=relative-heading-rank; base=h2

\section{第八十二章}

雷提丘斯，奥顿（在厄杜伊人中）的主教，在君士坦丁统治时期于高卢享有盛名。我读过他的\WorkTitle{雅歌注释}以及另一部巨著\WorkTitle{驳诺瓦提安}，但除此之外，未发现他的其他作品。\par

% structure: p0180 source=h2 target=section reason=relative-heading-rank; base=h2

\section{第八十三章}

美多迪乌斯，吕西亚的奥林普斯主教，后为提洛主教，撰写了文笔精炼、逻辑严密的\WorkTitle{驳波斐利}，以及\WorkTitle{十童女论宴}、优秀的著作\WorkTitle{论复活}（驳奥力振）、\WorkTitle{论交鬼的妇人}和\WorkTitle{论自由意志}（亦驳奥力振）。他还撰写了\WorkTitle{创世记注释}、\WorkTitle{雅歌注释}以及许多其他广泛流传的著作。在最近的迫害末期，或如其他人所说，在德西乌斯和瓦勒良统治时期，他在希腊的卡尔基斯殉道。\par

% structure: p0182 source=h2 target=section reason=relative-heading-rank; base=h2

\section{第八十四章}

尤文库斯，西班牙人，出身贵族，身为司铎，几乎逐字地将四福音书译成六音步诗行，撰写了四卷。他还以同样的格律，创作了涉及圣事次序的其他作品。他活跃于君士坦丁统治时期。\par

% structure: p0184 source=h2 target=section reason=relative-heading-rank; base=h2

\section{第八十五章}

欧斯塔提乌斯，潘菲利亚人，来自西德，先为叙利亚贝罗亚主教，后为安提约基雅主教，管理教会，撰写了大量驳斥亚略派学说的著作，在君士坦提乌斯皇帝统治下被流放至色雷斯的特拉亚诺波利斯，直至今日。他的著作现存有\WorkTitle{论灵魂}、\WorkTitle{论腹语术，驳奥力振}以及数量众多的\WorkTitle{书信}。\par

% structure: p0186 source=h2 target=section reason=relative-heading-rank; base=h2

\section{第八十六章}

安基拉的\Person{马尔塞鲁}主教，在\Person{君士坦丁}和\Person{君士坦提乌斯}在位时期活跃，撰写了多卷不同\WorkTitle{命题}，尤其驳斥亚略派。\Person{亚斯德留}和\Person{阿波利纳里斯}的著作驳斥他，指责他有撒贝略主义。\Person{希拉里}也在其著作\WorkTitle{驳斥亚略派}第七卷中提及他为异端者，但他为自己辩护，因为\Place{罗马}主教\Person{尤利乌}和\Place{亚历山大}的\Person{亚大纳削}曾与他共融。\par

% structure: p0188 source=h2 target=section reason=relative-heading-rank; base=h2

\section{第八十七章}

\Place{亚历山大}的主教\Person{亚大纳削}，因亚略派的诡计而受窘迫，逃往\Place{高卢}的\Person{君士坦斯}皇帝处。携信返回后，皇帝驾崩，再次逃亡躲藏，直至\Person{若维安}登基，他返回教会，并在\Person{瓦伦斯}在位时期去世。流传着他的多种著作：两卷\WorkTitle{驳斥外邦人}，一卷\WorkTitle{驳斥瓦伦斯与乌尔撒基乌斯}、\WorkTitle{论童贞}，多卷\WorkTitle{论亚略派的迫害}，还有\WorkTitle{论圣咏标题}和\WorkTitle{隐修士安当传}，以及\WorkTitle{庆节书信}和其他不胜枚举的著作。\par

% structure: p0190 source=h2 target=section reason=relative-heading-rank; base=h2

\section{第八十八章}

隐修士\Person{安当}，\Place{亚历山大}的主教\Person{亚大纳削}曾为其撰写长篇传记。他用科普特语向各隐修院发出七封书信，这些书信在思想和语言上具有宗徒风范，并被译成希腊文。其中最主要的一封是\WorkTitle{致阿尔塞诺特人}。他在\Person{君士坦丁}及其儿子们统治时期活跃。\par

% structure: p0192 source=h2 target=section reason=relative-heading-rank; base=h2

\section{第八十九章}

\Place{安基拉}的主教\Person{巴西略}，医学博士，撰写了\WorkTitle{驳斥马尔塞鲁及论童贞}及其他一些著作——在\Person{君士坦提乌斯}在位时期，他与\Place{塞巴斯特}的\Person{欧斯塔修}一同任\Place{马其顿}的首席主教。\par

% structure: p0194 source=h2 target=section reason=relative-heading-rank; base=h2

\section{第九十章}

\Place{色雷斯}\Place{赫拉克利亚}的主教\Person{德奥多若}，在\Person{君士坦提乌斯}皇帝在位时期出版了注释：\WorkTitle{论玛窦及若望}、\WorkTitle{论书信}和\WorkTitle{论圣咏}。这些注释文笔优美清晰，展现出卓越的历史感。\par

% structure: p0196 source=h2 target=section reason=relative-heading-rank; base=h2

\section{第九十一章}

\Place{厄梅萨}的\Person{欧瑟比}，拥有出色的修辞才华，撰写了无数适合赢得大众喝彩的著作，他以历史笔法写作，为从事公开演说者所勤读。其中主要作品有：\WorkTitle{驳斥犹太人、外邦人和诺瓦替安派}以及\WorkTitle{福音讲道集}，篇幅短小但数量众多。他在\Person{君士坦提乌斯}皇帝在位时期活跃并去世，葬于\Place{安提约基雅}。\par

% structure: p0198 source=h2 target=section reason=relative-heading-rank; base=h2

\section{第九十二章}

\Place{塞浦路斯}\Place{勒德拉}（或\Place{勒乌科特昂}）的主教\Person{特里斐利}，是他那个时代最雄辩的人，在\Person{君士坦提乌斯}在位时期声名卓著。我读过他的\WorkTitle{雅歌注释}。据说他还撰写了其他许多著作，但均未传至我们手中。\par

% structure: p0200 source=h2 target=section reason=relative-heading-rank; base=h2

\section{第九十三章}

\Person{多纳图斯}，多纳图派由此人在\Place{非洲}兴起于\Person{君士坦丁}和\Person{君士坦提乌斯}皇帝在位时期，他声称正统派在迫害期间将圣经交给了异教徒，并以其说服力几乎欺骗了全\Place{非洲}，尤其是\Place{努米底亚}。他的许多著作与他的异端有关，流传至今，包括\WorkTitle{论圣神}，这是一部在教义上属亚略派的作品。\par

% structure: p0202 source=h2 target=section reason=relative-heading-rank; base=h2

\section{第九十四章}

\Person{亚斯德留}，亚略派的哲学家，在\Person{君士坦提乌斯}在位时期撰写了注释：\WorkTitle{论罗马人书}、\WorkTitle{论福音}和\WorkTitle{论圣咏}，以及许多其他为他派中人所勤读的著作。\par

% structure: p0204 source=h2 target=section reason=relative-heading-rank; base=h2

\section{第九十五章}

\Place{卡利亚里}的主教\Person{路济斐尔}，由主教\Person{利贝里乌斯}与\Place{罗马}教会的司铎\Person{潘克拉修}和\Person{希拉里}一起派往\Person{君士坦提乌斯}皇帝处，作为信仰的使节。当他拒绝谴责以\Person{亚大纳削}为代表的尼西亚信仰时，被再次遣往\Place{巴勒斯坦}，他以非凡的坚毅和殉道精神，写了一本书驳斥\Person{君士坦提乌斯}皇帝，并送给他阅读；不久后，在\Person{尤利安}皇帝在位时期返回\Place{卡利亚里}，并在\Person{瓦伦提尼安}在位时期去世。\par

% structure: p0206 source=h2 target=section reason=relative-heading-rank; base=h2

\section{第九十六章}

\Person{欧瑟比}，\Place{撒丁}人，最初在\Place{罗马}担任读经员，后任\Place{韦尔切利}主教。因信仰宣认被\Person{君士坦提乌斯}皇帝遣往\Place{西托波利斯}，后又至\Place{卡帕多细亚}，在\Person{尤利安}皇帝在位时期返回教会，并出版了\Place{凯撒勒雅}的\Person{欧瑟比}的\WorkTitle{圣咏注释}，由希腊文译为拉丁文，在\Person{瓦伦提尼安}和\Person{瓦伦斯}在位时期去世。\par

% structure: p0208 source=h2 target=section reason=relative-heading-rank; base=h2

\section{第九十七章}

\Person{福图纳提安}，\Place{非洲}人，\Place{阿奎利亚}主教，在\Person{君士坦提乌斯}在位时期撰写了简短的\WorkTitle{福音注释}，按章节编排，文风质朴；他之所以受人憎恶，是因为当\Place{罗马}主教\Person{利贝里乌斯}因信仰被流放时，\Person{利贝里乌斯}受\Person{福图纳提安}的催促所诱导，签附了异端。\par

% structure: p0210 source=h2 target=section reason=relative-heading-rank; base=h2

\section{第九十八章}

\Person{阿卡基乌斯}，因一目失明，被人戏称为\Quote{独眼}，\Place{巴勒斯坦}\Place{凯撒勒雅}教会的主教，撰写了十七卷\WorkTitle{训道篇}和六卷\WorkTitle{杂问}，以及许多其他论题的文章。他在\Person{君士坦提乌斯}皇帝在位时期势力如此之大，以致他使\Person{菲利克斯}取代\Person{利贝里乌斯}成为\Place{罗马}主教。\par

% structure: p0212 source=h2 target=section reason=relative-heading-rank; base=h2

\section{第九十九章}

\Place{特穆依斯}的主教\Person{塞拉皮翁}，因其卓越才华被誉为\Quote{经师}，是隐修士\Person{安当}的密友。他出版了一部优秀的著作\WorkTitle{驳斥摩尼教徒}，以及另一部\WorkTitle{论圣咏标题}，还有致不同人士的宝贵\WorkTitle{书信}。在\Person{君士坦提乌斯}皇帝在位时期，他作为精修圣人而闻名。\par

% structure: p0214 source=h2 target=section reason=relative-heading-rank; base=h2

\section{第一百章}

\Person{希拉略}，阿基坦的\Place{普瓦捷}主教，曾是阿尔勒主教\Person{撒图尔尼努斯}的同党。他被\Place{贝济耶}会议流放到\Place{弗里吉亚}，期间撰写了十二卷《\WorkTitle{反亚略派}》、另一部写给高卢主教们的《\WorkTitle{论议会}》，以及《\WorkTitle{圣咏注释}》，即对第一首和第二首、从第五十一首到第六十二首、以及从第一百一十八首到卷末的注释。在此著作中，他模仿了\Person{奥利振}，但也加入了一些独创内容。他有一本小册子《\WorkTitle{致君士坦提乌斯}》，是在他居留\Place{君士坦丁堡}时呈献给皇帝的；另一本《\WorkTitle{论君士坦提乌斯}》是他死后所写；还有一本《\WorkTitle{反瓦伦斯和乌尔撒基乌}》，其中包含阿里米尼会议和塞琉西亚会议的历史；以及《\WorkTitle{致总督撒路斯特}》或《\WorkTitle{反狄奥斯库鲁}》；还有一部《\WorkTitle{赞美诗与奥迹}》，一部《\WorkTitle{玛窦福音注释}》，以及关于《\WorkTitle{约伯传注释}》的论文（他根据\Person{奥利振}的希腊文本自由翻译）；另一部精巧的小作品《\WorkTitle{反奥克森提乌}》以及致不同人士的《\WorkTitle{书信集}》。据说他还写了《\WorkTitle{雅歌注释}》，但这部著作我们并不知道。他在\Person{瓦伦提尼安}和\Person{瓦伦斯}在位期间卒于\Place{普瓦捷}。\par

% structure: p0216 source=h2 target=section reason=relative-heading-rank; base=h2

\section{第一百零一章}

\Person{维克多利努斯}，非洲裔，在\Person{君士坦提乌斯}皇帝治下于\Place{罗马}教授修辞学；在极晚年，他皈依\Person{基督}信仰，撰写了反对\Person{亚略}的著作，采用辩证文体，语言极为晦涩，只有学者才能理解。他还撰写了《\WorkTitle{书信注释}》。\par

% structure: p0218 source=h2 target=section reason=relative-heading-rank; base=h2

\section{第一百零二章}

\Person{提督}，\Place{博斯特拉}主教，在皇帝\Person{尤利安}和\Person{约维尼安}统治期间撰写了反对摩尼教的激烈著作，以及其他一些作品。他在\Person{瓦伦斯}治下去世。\par

% structure: p0220 source=h2 target=section reason=relative-heading-rank; base=h2

\section{第一百零三章}

\Person{达玛苏斯}，\Place{罗马}主教，具有杰出的作诗天赋，发表了许多英雄格律的短篇作品。他在\Person{德奥多西}皇帝统治期间去世，享年将近八十岁。\par

% structure: p0222 source=h2 target=section reason=relative-heading-rank; base=h2

\section{第一百零四章}

\Person{阿波利纳利斯}，\Place{叙利亚}地区\Place{劳迪西亚}主教，是一位长老的儿子。年轻时勤奋学习文法，后来撰写了无数卷关于圣经的著作，在\Person{德奥多西}皇帝统治期间去世。现存有他三十卷《\WorkTitle{反波菲利}》，通常被视为其最佳著作之一。\par

% structure: p0224 source=h2 target=section reason=relative-heading-rank; base=h2

\section{第一百零五章}

\Person{额我略}，\Place{贝提卡}地区\Place{埃尔维拉}主教，写作直至极晚年，以平庸的文笔撰写了多种论文，以及一部文笔优雅的《\WorkTitle{论信德}》。据说他仍在世。\par

% structure: p0226 source=h2 target=section reason=relative-heading-rank; base=h2

\section{第一百零六章}

\Person{帕西亚努斯}，\Place{比利牛斯山脉}地区的\Place{巴塞罗那}主教，一位贞洁而雄辩的人，其生平与言辞同样卓越。他撰写了多种短篇作品，其中包括《\WorkTitle{鹿}》和《\WorkTitle{反诺瓦提安派}》。在皇帝\Person{德奥多西}统治期间死于极晚年。\par

% structure: p0228 source=h2 target=section reason=relative-heading-rank; base=h2

\section{第一百零七章}

\Person{福提努斯}，\Place{加拉提亚}人，\Person{马尔塞鲁}的门徒，被按立为\Place{西尔米乌姆}主教。他试图引入厄便尼异端，后来被皇帝\Person{瓦伦提尼安}逐出教会，期间撰写了多卷著作，其中最出色的是《\WorkTitle{反外邦}》和《\WorkTitle{致瓦伦提尼安}》。\par

% structure: p0230 source=h2 target=section reason=relative-heading-rank; base=h2

\section{第一百零八章}

\Person{费巴狄乌斯}，\Place{高卢}地区\Place{阿让}主教，出版了一部《\WorkTitle{反亚略派}》。据称他还有其他著作，但我尚未读到。他仍健在，但因年老而虚弱。\par

% structure: p0232 source=h2 target=section reason=relative-heading-rank; base=h2

\section{第一百零九章}

\Person{狄迪慕}，\Place{亚历山大}人，幼年失明，因而未受启蒙教育，但他展现出如此惊人的智力，以至于完美掌握了辩证术甚至几何学（这些科学尤其需要视力）。他撰写了众多令人钦佩的著作：《\WorkTitle{全部圣咏注释}》、《\WorkTitle{玛窦及若望福音注释}》、《\WorkTitle{论教义}》，还有两部《\WorkTitle{反亚略派}》，一部《\WorkTitle{论圣神}》（我将其译成了拉丁文），十八卷《\WorkTitle{依撒意亚注释}》，三卷《\WorkTitle{欧瑟亚注释}》（致我），五卷《\WorkTitle{匝加利亚注释}》（应我之请而写），以及《\WorkTitle{约伯传注释}》等许多其他作品，若要一一列举，本身便可成为一部著作。他仍健在，已超过八十三岁。\par

% structure: p0234 source=h2 target=section reason=relative-heading-rank; base=h2

\section{第一百一十章}

\Person{奥普塔图斯}，非洲人，\Place{米勒维斯}主教，在\Person{瓦伦提尼安}和\Person{瓦伦斯}皇帝统治期间，代表天主教派撰写了六卷书反对多纳图派的诽谤，在其中他断言多纳图派的罪行被错误地归咎于天主教派。\par

% structure: p0236 source=h2 target=section reason=relative-heading-rank; base=h2

\section{第一百一十一章}

\Person{阿基略·塞维鲁}，西班牙人，属于那位\Person{塞维鲁}的家族，\Person{拉克坦提乌斯}的两卷《\WorkTitle{书信集}》就是致他的。他撰写了一卷诗文混杂的作品，可视为其全部生活的指南，他称之为《\WorkTitle{灾难或考验}》。他在\Person{瓦伦提尼安}统治期间去世。\par

% structure: p0238 source=h2 target=section reason=relative-heading-rank; base=h2

\section{第一百一十二章}

\Person{济利禄}，\Place{耶路撒冷}主教，多次被教会驱逐，最后又被接纳。他在\Person{德奥多西}统治期间连续担任主教职位八年。他年轻时编纂的某些《\WorkTitle{教理讲授}》今日仍存。\par

% structure: p0240 source=h2 target=section reason=relative-heading-rank; base=h2

\section{第一百一十三章}

\Person{欧佐伊乌斯}年轻时，与\Place{纳齐安}主教\Person{额我略}一起，在\Place{凯撒利亚}跟随修辞学家\Person{特斯佩西乌斯}受教。后来当他担任同一城市的主教时，不辞辛劳试图修复由\Person{奥利振}和\Person{庞菲禄}收集、已遭损毁的图书馆。最后，在皇帝\Person{德奥多西}统治期间，他被逐出教会。他的多种论文流传于世，人们很容易读到它们。\par

% structure: p0242 source=h2 target=section reason=relative-heading-rank; base=h2

\section{第一百一十四章}

\Person{埃皮法尼乌斯}，\Place{塞浦路斯}岛\Place{萨拉米纳}主教，撰写了《\WorkTitle{反一切异端}》以及其他许多著作，学者因其主题而争先阅读，平民则因其语言而喜爱。他仍健在，在极晚年创作各种短篇作品。\par

% structure: p0244 source=h2 target=section reason=relative-heading-rank; base=h2

\section{第一百一十五章}

\Person{厄弗冷}，\Place{厄德萨}教会执事，用叙利亚语撰写了众多著作，他声名卓著，以致在某些教会中，在宣读圣经之后，会公开诵读他的作品。\par

我曾读过他一部希腊文著作《\WorkTitle{论圣神}》，是有人从叙利亚语翻译过来的；即使在译文中，我也能认出那崇高天才的敏锐力量。他在\Person{瓦伦斯}统治期间去世。\par

% structure: p0247 source=h2 target=section reason=relative-heading-rank; base=h2

\section{第一百一十六章}

卡帕多细亚凯撒勒雅（昔称马扎卡）的主教\Person{巴西略}，撰写了精心之作\WorkTitle{驳欧诺弥}、一卷\WorkTitle{论圣神}、九篇\WorkTitle{论六日创造}讲道辞，以及一部\WorkTitle{论苦修}和若干短篇论著。他在格拉提安皇帝在位期间去世。\par

% structure: p0249 source=h2 target=section reason=relative-heading-rank; base=h2

\section{第117章}

\Person{额我略}，纳齐盎的主教，是一位极富口才的人，也是我研读圣经的导师。他的作品总计约三万行，其中包括\WorkTitle{论其弟凯撒留之死}、\WorkTitle{论爱德}、\WorkTitle{赞玛加伯}、\WorkTitle{赞西彼廉}、\WorkTitle{赞阿塔纳修}、\WorkTitle{赞哲学家玛克西默——在他流放归来之后}。然而，后者有人以赫罗纳的假名题写，因为额我略另有作品谴责同一位玛克西默，仿佛依据场合需要，可以有时称赞、有时谴责同一个人。他的其他作品还有：一篇六音步诗体作品，题为\WorkTitle{论贞洁与婚姻之辩}；两卷\WorkTitle{驳欧诺弥}；一卷\WorkTitle{论圣神}；以及一卷\WorkTitle{驳皇帝尤利安}。他的演说风格追随波勒蒙。他在有生之年指定了主教继任者，然后退隐乡间，过隐修士的生活，并于三年前或更早，在德奥多西皇帝在位期间去世。\par

% structure: p0251 source=h2 target=section reason=relative-heading-rank; base=h2

\section{第118章}

\Person{路济乌}，亚流派人士，继阿塔纳修之后担任亚历山大教会的主教，直至德奥多西皇帝在位期间被罢黜。他的若干\WorkTitle{论逾越节}庆节书信以及少量\WorkTitle{杂论}短篇作品流传于世。\par

% structure: p0253 source=h2 target=section reason=relative-heading-rank; base=h2

\section{第119章}

塔尔索的主教\Person{狄奥多若}，在仍担任安提约基雅长老时便享有盛名。他的\WorkTitle{论宗徒书信释义}以及许多其他作品流传于世，其风格仿效伟大的厄梅森人欧瑟比，虽追随其思想，但因不通世俗文学，未能模仿其文采。\par

% structure: p0255 source=h2 target=section reason=relative-heading-rank; base=h2

\section{第120章}

基齐库的主教\Person{欧诺弥}，亚流派人士，在其异端中公然亵渎，公开宣扬他人隐晦的言论。据说他仍住在卡帕多细亚，并写作大量反对教会的作品。\Person{阿波里纳}、\Person{狄迪慕}、凯撒勒雅的\Person{巴西略}、纳齐盎的\Person{额我略}以及尼撒的\Person{额我略}都曾撰文答复他。\par

% structure: p0257 source=h2 target=section reason=relative-heading-rank; base=h2

\section{第121章}

阿比拉的主教\Person{普里西利安}，属希达修与依达修派，在特里尔被僭主玛克西默处死。他发表了多篇短论，其中有些已传至我们手中。至今仍有人指控他沾染诺斯替主义，即依肋内所记述的巴西里德或玛尔谷的异端；而他的辩护者则声称他绝无此等思想。\par

% structure: p0259 source=h2 target=section reason=relative-heading-rank; base=h2

\section{第122章}

西班牙的\Person{拉特洛尼安}，学识渊博，在诗作方面堪与古代诗人媲美，也随\Person{普里西利安}、\Person{斐利奇西默}、\Person{犹利安}和\Person{欧克罗提亚}一同在特里尔被处死，他们是分裂的共谋者。他用不同诗体写成的各种才华之作流传于世。\par

% structure: p0261 source=h2 target=section reason=relative-heading-rank; base=h2

\section{第123章}

贝提卡的\Person{提贝里安}，为回应有人暗示他附和普里西利安的异端，写了一篇华丽的杂体辩护辞。但在他的友人死后，他因流放的漫长而感到厌倦，改变了心意，正如圣经所载：\Quote{狗转回去吃自己所吐的}，并与一位献身于基督的童贞女（修女）结婚。\par

% structure: p0263 source=h2 target=section reason=relative-heading-rank; base=h2

\section{第124章}

米兰的主教\Person{盎博罗削}，目前仍在写作。我对他不作评判，因他尚在人世，惟恐要么称赞（沦为谄媚），要么指责（说了实话）。\par

% structure: p0265 source=h2 target=section reason=relative-heading-rank; base=h2

\section{第125章}

安提约基雅的主教\Person{厄瓦格利}，思想极为敏锐。他仍在担任长老时，便将尚未发表的各类论著读给我听。他还将阿塔纳修希腊文所著的\WorkTitle{真福安当生平}译成了我们的语言。\par

% structure: p0267 source=h2 target=section reason=relative-heading-rank; base=h2

\section{第126章}

亚历山大的\Person{盎博罗削}，\Person{狄迪慕}的门徒，撰写了长篇\WorkTitle{论教义}以驳斥阿波里纳，并且据最近有人告知我，还写了\WorkTitle{\BibleBook{约伯}{传}释义}。他仍在世。\par

% structure: p0269 source=h2 target=section reason=relative-heading-rank; base=h2

\section{第127章}

哲学家\Person{玛克西默}，生于亚历山大，在君士坦丁堡被祝圣为主教，后被罢黜。他写了一部卓越的\WorkTitle{论信德}驳斥亚流派，并在米兰呈献给格拉提安皇帝。\par

% structure: p0271 source=h2 target=section reason=relative-heading-rank; base=h2

\section{第128章}

尼撒的主教\Person{额我略}，凯撒勒雅的巴西略的弟弟，数年前曾将其驳欧诺弥的作品读给纳齐盎的额我略和我听。据说他还写了其他许多作品，并仍在写作。\par

% structure: p0273 source=h2 target=section reason=relative-heading-rank; base=h2

\section{第129章}

安提约基雅教会的长老\Person{若望}，厄梅森的欧瑟比与狄奥多若的追随者，据说撰写了许多书，但其中我只读过他的\WorkTitle{论司祭职}。\par

% structure: p0275 source=h2 target=section reason=relative-heading-rank; base=h2

\section{第130章}

巴勒斯坦凯撒勒雅的主教\Person{格拉修}，继欧佐依之后，据说他或多或少以精雕细琢的风格写作，但未发表其作品。\par

% structure: p0277 source=h2 target=section reason=relative-heading-rank; base=h2

\section{第131章}

斯基提亚托米的主教\Person{德奥提慕}，发表了简短而警句式的对话体论著，风格古旧。我听说他正在写其他作品。\par

% structure: p0279 source=h2 target=section reason=relative-heading-rank; base=h2

\section{第132章}

\Person{德克斯德}，我前面提及的\Person{帕西安}之子，在其时代出类拔萃，笃信基督信仰。据说他撰写了\WorkTitle{普世史}，但我尚未读过。\par

% structure: p0281 source=h2 target=section reason=relative-heading-rank; base=h2

\section{第133章}

依科尼雍的主教\Person{安斐罗基}，最近读给我听了一卷\WorkTitle{论圣神}，论证圣神是天主，应受敬拜，且是全能的。\par

% structure: p0283 source=h2 target=section reason=relative-heading-rank; base=h2

\section{第134章}

\Person{索弗罗尼}，学识超群。他年少时写了\WorkTitle{白冷赞}，近来有一部重要作品\WorkTitle{论塞拉皮斯之倾覆}，以及致\Person{欧斯塔基}的\WorkTitle{论童贞}和\WorkTitle{隐修士希拉略传}。他将我的短篇作品以极精良的风格译成希腊文，也将我从希伯来文译为拉丁文的\WorkTitle{\BibleBook{圣咏}{集}}和\WorkTitle{\BibleBook{}{先知书}}译成了希腊文。\par

% structure: p0285 source=h2 target=section reason=relative-heading-rank; base=h2

\section{第135章}

我，\Person{耶柔米}，\Person{厄乌瑟比}之子，\Place{斯特里多}城人——该城位于\Place{达尔马提亚}与\Place{潘诺尼亚}交界，曾被\Person{哥特人}摧毁——时至今年，即\Person{狄奥多西}皇帝在位第十四年，已写成下列作品：\WorkTitle{隐修士保禄传}一卷，\WorkTitle{致不同人士书信}一卷，\WorkTitle{劝勉\Person{赫略多若}}，\WorkTitle{\Person{路济斐里亚诺}与正统派之争辩}，\WorkTitle{普世编年史}，\WorkTitle{\Person{奥力振}论\Person{耶肋米亚}与\Person{厄则克耳}讲道集二十八篇}，我从希腊文译成拉丁文；\WorkTitle{论色辣芬}，\WorkTitle{论贺三纳}，\WorkTitle{论谨慎与浪子}，\WorkTitle{论古律三题}，\WorkTitle{雅歌讲道集}两篇，\WorkTitle{反\Person{赫耳维狄}}，\WorkTitle{论\Person{玛利亚}终身童贞}，\WorkTitle{致\Person{欧斯托基}论守童贞}一卷，\WorkTitle{致\Person{玛尔克拉}书信集}一卷，\WorkTitle{致\Person{保拉}慰唁信：论女儿之死}一卷，\WorkTitle{\Person{保禄}致\BibleBook{迦拉达}{书}注释}三卷，\WorkTitle{致\BibleBook{厄弗所}{书}注释}三卷，\WorkTitle{致\BibleBook{弟铎}{书}注释}一卷，\WorkTitle{致\BibleBook{费肋孟}{书}注释}一卷，\WorkTitle{\BibleBook{训道}{篇}注释}，\WorkTitle{\BibleBook{}{创世纪}希伯来问题}一卷，\WorkTitle{\Place{犹太}地名录}一卷，\WorkTitle{希伯来名称}一卷，\WorkTitle{\Person{狄迪莫斯}论圣神}——我从希腊文译成拉丁文——一卷，\WorkTitle{\BibleBook{路加}{福音}讲道集三十九篇}，\WorkTitle{\BibleBook{}{圣咏集}第十至十六篇}七卷，\WorkTitle{论被掳隐修士}，\WorkTitle{真福\Person{希拉良}传}。我将\WorkTitle{新约}从希腊文译出，将\WorkTitle{旧约}从希伯来文译出，而\WorkTitle{致\Person{保拉}与\Person{欧斯托基}书信}我写了多少已不清楚，因我几乎每日写作。此外，我还写了\WorkTitle{\BibleBook{}{米该亚}释义}两卷，\WorkTitle{\BibleBook{}{纳鸿}释义}一卷，\WorkTitle{\BibleBook{}{哈巴谷}释义}两卷，\WorkTitle{\BibleBook{}{索福尼亚}释义}一卷，\WorkTitle{\BibleBook{}{哈盖}释义}一卷，以及许多其它论先知的作品，尚未完成，仍在继续撰写中。\par

\SourceRecord{Translated by Ernest Cushing Richardson.；Nicene and Post-Nicene Fathers, Second Series；Vol. 3.；Edited by Philip Schaff and Henry Wace.；Buffalo, NY: Christian Literature Publishing Co.,；1892.；<http://www.newadvent.org/fathers/2708.htm>.}
